\documentclass[11pt]{ctexart}
\usepackage[a4paper,margin=1in]{geometry}
\usepackage{graphicx}
\usepackage{xcolor}
\usepackage{hyperref}
\definecolor{refgray}{gray}{0.45}
\title{\textbf{Market Views｜全球投行叙事汇编}\\[0.4em]{\large 宏观主线从总量波动转向结构性分化，AI投资从算力军备转向效率优先，大宗商品从短缺溢价转向供给修复与政策博弈}}
\author{KC桌面 自动生成}
\date{260711}
\begin{document}
\maketitle
\tableofcontents
\newpage
\section{一页摘要}
\begin{itemize}
  \item 美国中期选举后政策不确定性回落驱动市场风格切换，科技、成长、质量风格领跑，标普500通常录得超额回报。
  \item 亚洲厄尔尼诺对农业产出冲击显著（平均-2.4\%），但油价温和、灌溉改善及政府干预缓冲食品通胀，需求侧冲击（农村消费下降）才是被低估的风险。
  \item 韩元在强劲贸易顺差下持续走弱，根源在于结构性资金外流与外资对冲行为切断了贸易盈余向汇率的传导，传统定价框架失效。
  \item 全球铝市场2027年将从短缺转为明显过剩，中东产能恢复速度远超预期，印尼新增产能加速投放，市场此前普遍预期短缺将持续。
  \item AI投资正从算力军备转向效率优先，记忆体长协、HDD、硅片因AI需求出现结构性量价齐升，中国开源模型逼近‘足够好’临界点。
  \item 人形机器人量产节奏快于预期，2026年Q3-Q4产量将从每周约100台爬坡至约1.9万台，中国供应链企业收入占比快速提升。
  \item 新兴市场指数成分已发生根本性变化，前三大重仓股（台积电、三星、SK海力士）占比近30\%，指数集中度超过标普500。
  \item 美国固废行业RNG业务EBITDA利润率高达40-50\%，但90\%营收依赖D3 RIN信用额度，信用价格波动（2-3.75美元区间）将导致EBITDA从5.13亿至10.26亿美元的巨幅变化。
\end{itemize}
\section{全球投行叙事汇编}
今日纳入 59 篇投行/券商研究，按 15 个市场、资产与产业主题系统整理。
\newpage
\subsection{宏观与政策：不确定性节奏与结构性分化}
\textbf{全球宏观的核心矛盾正从总量波动转向结构性分化：美国中期选举后政策不确定性回落驱动市场风格切换，亚洲厄尔尼诺冲击需求侧而非通胀，中国通胀呈现上游热下游冷格局，日本央行独立性边界成为市场焦点，人民币国际化则进入管道式开放新阶段。}
\subsubsection*{主流共识}
\begin{itemize}
  \item 美国中期选举后12个月标普500通常录得超额回报，科技、成长、质量风格领跑，核心驱动力是政策不确定性从高位回落。
  \item 亚洲厄尔尼诺对农业产出冲击显著（平均-2.4\%），但食品通胀未必同步上升，因油价温和、灌溉改善及政府干预形成缓冲。
  \item 中国通胀读数由外部输入主导，核心CPI约0.4\%低位稳定，内需定价领域（猪肉、下游消费品）持续偏弱。
\end{itemize}
\subsubsection*{分歧与条件差异}
\begin{itemize}
  \item 厄尔尼诺对亚洲经济的影响路径存在分歧：市场普遍关注食品通胀与央行政策，但NOM认为需求侧冲击（农村消费下降）才是真正被低估的风险。
  \item 日本基本政策修订中措辞调整是否足以消除市场对央行独立性的担忧：DB认为残留表述（“不重回通缩”“工资物价良性循环”）仍构成隐形门槛，市场并未放松警惕。
  \item 人民币国际化新框架的开放程度：MS判断为“不对称开放”——收紧非正式资本流出，同时扩大香港离岸市场正规渠道，并非全面开闸。
\end{itemize}
\subsubsection*{机构观点}
\paragraph{BARC} 美国中期选举后市场超额回报的核心驱动力并非宏观基本面改善，而是政策不确定性（EPU）从高位快速回落。历史数据显示，选举前8-9月标普500弱势明显，但选举后12个月科技、成长、质量风格显著领跑。基准情景为分裂政府（共和党总统、民主党众议院、共和党参议院），该格局下政策不确定性处于历史低位。
{\small\color{refgray}数据：中期选举后12个月标普500超额回报显著高于非选举年份均值\par}
{\small\color{refgray}数据：选举前8-9月标普500与非选举年份的回报差异最大\par}
{\small\color{refgray}数据：科技股超额收益与EPU节奏变化正相关，与EPU绝对水平无关\par}
{\small\color{refgray}边际变化：市场关注点从选举结果本身转向政策不确定性下降的节奏与幅度，若不确定性未如期回落，市场逻辑将动摇。\par}
\paragraph{NOM} 厄尔尼诺对亚洲经济的真正冲击不在通胀，而在需求侧。过去25年七次强厄尔尼诺期间，亚洲农业产出平均下降2.4个百分点，但食品通胀反而意外下降，因油价温和、灌溉改善及政府干预。菲律宾和印度暴露度最高：菲律宾为净食品进口国，大米CPI权重8.9\%；印度夏季作物播种面积同比下降20.8\%。央行大概率不因供给冲击调整利率，但贸易保护主义（出口限制、关税）是最大潜在风险。
{\small\color{refgray}数据：过去25年七次厄尔尼诺，亚洲农业产出平均下降2.4个百分点\par}
{\small\color{refgray}数据：菲律宾大米在CPI中权重8.9\%，食品进口占GDP约2\%\par}
{\small\color{refgray}数据：印度夏季作物播种面积同比下降20.8\%，其中水稻-13.1\%，油籽-39.3\%\par}
{\small\color{refgray}数据：泰国米价+16.5\%，豆油+58\%\par}
{\small\color{refgray}边际变化：市场焦点从食品通胀转向需求侧冲击（农村收入下降影响消费）及贸易保护主义风险，央行政策预期从加息转向观望。\par}
\paragraph{NOM} 中国6月通胀数据呈现“上游热、下游冷”格局：PPI同比4.1\%但涨幅集中于石油（贡献1.6pp）和有色金属（贡献1.7pp），下游消费品PPI同比-0.9\%且持续负值；CPI同比1.0\%回落，0.23pp降幅几乎全部来自石油和黄金贡献减少，剔除后核心CPI约0.4\%低位稳定。内需驱动的涨价动力未形成，企业利润空间被两头挤压。
{\small\color{refgray}数据：6月CPI同比1.0\%（前值1.2\%），低于预期\par}
{\small\color{refgray}数据：原油和黄金对CPI拉动从0.83pp降至0.60pp\par}
{\small\color{refgray}数据：剔除原油和黄金后核心CPI约0.4\%\par}
{\small\color{refgray}数据：PPI同比4.1\%，石油+有色金属合计贡献超80\%\par}
{\small\color{refgray}边际变化：通胀读数由外部输入主导转向内需定价领域持续偏弱，市场对货币政策宽松预期升温。\par}
\paragraph{DB} 日本基本政策修订中货币政策措辞调整（从“实现强经济”改为“有助于稳定物价上涨”）是沟通上的积极一步，但不足以消除债券市场警惕。残留表述“在不重返通缩前提下实现物价稳定”和“确认工资与物价的良性循环”可能暗示政府对央行独立性施加了隐形门槛，这些标准在《日本银行法》和联合声明中均不存在。当日日本国债收益率曲线全面上行（30年期除外），市场并未放松戒备。
{\small\color{refgray}数据：修订后措辞明确央行首要目标为物价稳定\par}
{\small\color{refgray}数据：残留表述“不重返通缩”和“工资物价良性循环”在现行法律中无依据\par}
{\small\color{refgray}数据：修订当日日本国债收益率曲线除30年期外全面上行\par}
{\small\color{refgray}边际变化：市场关注点从措辞调整转向政府是否真正尊重央行独立性边界，政策沟通的边际改善被残留表述的不确定性抵消。\par}
\paragraph{MS} 人民币国际化进入管道式开放新阶段，核心是扩大香港离岸人民币市场深度与广度。新措施针对三大瓶颈：南向债券通额度从5000亿升至8000亿，香港人民币资金池从2000亿扩至5000亿且期限延至3年，推出5年期离岸人民币国债期货。同时推动人民币在大宗商品定价功能（跨境黄金清算、商品期货）。整体思路为“不对称开放”——收紧非正式资本流出，扩大正规渠道。
{\small\color{refgray}数据：南向债券通额度从5000亿升至8000亿人民币\par}
{\small\color{refgray}数据：香港人民币资金池从2000亿扩至5000亿，期限延至3年\par}
{\small\color{refgray}数据：首次推出5年期离岸人民币国债期货\par}
{\small\color{refgray}数据：推动跨境黄金清算结算机制\par}
{\small\color{refgray}边际变化：政策重心从堵漏转向建设性管道开放，离岸市场从“蓄水池”升级为可定价、可对冲、可深度配置的市场。\par}
\subsubsection*{关键数据}
\begin{itemize}
  \item 美国中期选举后12个月标普500超额回报显著高于非选举年份均值
  \item 亚洲强厄尔尼诺期间农业产出平均下降2.4个百分点
  \item 菲律宾大米CPI权重8.9\%，印度夏季作物播种面积同比-20.8\%
  \item 中国6月核心CPI约0.4\%，下游消费品PPI同比-0.9\%
  \item 日本基本政策修订残留“不重返通缩”表述，市场警惕未消
  \item 南向债券通额度升至8000亿，香港人民币资金池扩至5000亿
\end{itemize}
\begin{figure}[htbp]
\centering
\includegraphics[width=0.88\linewidth]{figures/fig_006.jpg}
\caption{FIGURE 2 - \#\# Equity returns pre- and post-midterm elections The "midterm effect" on U.S. equity returns is a well-documented phenomenon. Our prior research established that the S\&P 500 tends to deliver materially lower returns in the first 3 quarters of a U.S. midterm e}
\end{figure}
\begin{figure}[htbp]
\centering
\includegraphics[width=0.88\linewidth]{figures/fig_036.jpg}
\caption{Figure 2 - \#\# Introduction El Niño has officially formed, and is projected to be a strong-to-very-strong episode, with the potential to morph into a "super El Niño" later this year. For most Asian economies, El Niño results in below-normal rainfall and heatwaves. This ca}
\end{figure}
\begin{figure}[htbp]
\centering
\includegraphics[width=0.88\linewidth]{figures/fig_008.jpg}
\caption{FIGURE 3 - FIGURE 3. ...with the YTD return trajectory reshaped most meaningfully in the summer months FIGURE 4. Post-midterms, the S\&P 500 typically recovers the previous year's underperformance... FIGURE 5. ...by closing the gap meaningfully over the first 3 quarters o}
\end{figure}
\begin{figure}[htbp]
\centering
\includegraphics[width=0.88\linewidth]{figures/fig_038.jpg}
\caption{Figure 4 - \#\# Impact so far in Asia Fig. 4: India: Kharif (summer) crop area sown as of 5 July Fig. 5: Change in various food and fertilizer prices in 2026 \#\# Past El Niño cycles}
\end{figure}
{\small\color{refgray}本节综合 5 篇研究；涉及 BARC、NOM、DB、MS。}
\newpage
\subsection{亚洲汇率的结构性错位：从韩元弱势到人民币定价微调}
\textbf{传统汇率定价框架在亚洲多个市场失效：韩元在强劲贸易顺差下持续走弱，根源在于结构性资金外流与外资对冲行为切断了贸易盈余向汇率的传导；人民币中间价定价机制中逆周期因子使用力度微弱，显示政策层对汇率波动的容忍度上升；而以色列谢克尔与韩元在科技周期中的分化，则揭示了内资回流与外资流出两种截然不同的资金传导机制。}
\subsubsection*{主流共识}
\begin{itemize}
  \item 韩元弱势并非基本面问题，而是资金流动的结构性错位——企业未将美元结汇回流、外资增加对冲、个人配置海外资产，共同切断了贸易顺差向汇率传导的传统路径。
  \item 人民币中间价定价基本跟随市场收盘，逆周期因子影响极小，显示当前政策层对汇率走势的干预意愿不强。
  \item 科技周期对汇率的影响取决于资金传导机制：以色列谢克尔受益于内资回流，韩元则受外资流出压制，两者走势分化。
\end{itemize}
\subsubsection*{分歧与条件差异}
\begin{itemize}
  \item 韩元短期走势：MS认为若无美元整体走弱，USD/KRW可能维持在1490-1560区间；但DB指出科技板块调整可能通过ILS/KRW交叉盘放大韩元波动。
  \item 人民币中间价定价中逆周期因子的启用程度：NOM模型显示当前因子影响仅0.6个基点，但若后续因子调整幅度扩大或方向逆转，可能意味着政策态度变化。
  \item 韩元与Kospi的关联性：传统正相关已发生变化，MS和DB均观察到Kospi上涨时外资反而减持或增加对冲，但两者对本土资金能否支撑韩元存在不同判断。
\end{itemize}
\subsubsection*{机构观点}
\paragraph{MS} 韩元持续走弱的核心原因在于结构性资金外流和出口企业外汇转换滞后，切断了贸易盈余向汇率传导的传统路径。韩国经常账户盈余在2026年1-5月已达1410亿美元，但企业未将美元结汇回流，反而加大了对美直接投资（2025年达250亿美元，创历史新高）。同时，外资在配置韩国股票时越来越多地做外汇对冲，削弱了股汇联动。短期看，7-8月有半导体公司海外上市资金回流等支撑，但整体外资流出压力仍在，若无美元整体走弱，USD/KRW可能维持在1490-1560区间。
{\small\color{refgray}数据：韩国经常账户盈余2026年1-5月达1410亿美元\par}
{\small\color{refgray}数据：2025年韩国企业对美直接投资250亿美元，创历史新高\par}
{\small\color{refgray}数据：USD/KRW预计区间1490-1560\par}
{\small\color{refgray}边际变化：传统上经常账户盈余应推动本币升值，但MS指出资金流出结构（企业海外投资、外资对冲、个人海外配置）已使这一传导机制失效，这是对旧叙事（基本面改善→韩元走强）的颠覆。\par}
\paragraph{NOM} 人民币中间价定价模型显示，无逆周期因子的预测值（6.7936）较前一日官方收盘价仅高1个基点，计入逆周期因子后预测值（6.7942）较前一日定盘价低94个基点，但因子本身影响仅0.6个基点，说明当前政策层对汇率走势的干预意愿不强。模型误差连续多日较小，暗示定价机制基本跟随市场。下半年关键事件包括7月底政治局会议、10月国庆黄金周、11月深圳APEC会议、12月中央经济工作会议及可能的中美元首会晤，这些都可能成为汇率波动催化剂。
{\small\color{refgray}数据：无逆周期因子模型预测USD/CNY中间价6.7936\par}
{\small\color{refgray}数据：计入逆周期因子后预测值6.7942，较前一日定盘价低94个基点\par}
{\small\color{refgray}数据：逆周期因子影响仅0.6个基点\par}
{\small\color{refgray}边际变化：此前市场关注中间价点位本身，NOM模型强调逆周期因子的启用程度比点位更值得追踪——当前因子影响极小，若后续调整幅度扩大或方向逆转，将反映政策态度变化。\par}
\paragraph{DB} 以色列谢克尔与韩元在科技周期中走势分化，根源在于截然不同的资金传导机制。自2024年中期以来，ILS/KRW已走强约50\%，与超大规模云服务商涨幅同步。以色列方面，科技股上涨带动本土企业利润汇回和国内机构再平衡（卖出美元、买入谢克尔），形成内资回流支撑；韩国方面，Kospi上涨反而伴随外资减持或增加对冲，外资流出压过本土资金支撑。若科技板块出现明显调整，ILS/KRW交叉盘可能面临较大波动。
{\small\color{refgray}数据：ILS/KRW自2024年中期以来走强约50\%\par}
{\small\color{refgray}数据：以色列科技出口企业利润汇回驱动谢克尔走强\par}
{\small\color{refgray}数据：韩国Kospi上涨时外资倾向于减持或增加对冲\par}
{\small\color{refgray}边际变化：传统观点认为科技周期对类似经济体的货币影响趋同，DB拆解出内资回流（以色列）与外资流出（韩国）两种机制，解释了ILS/KRW的持续分化。\par}
\subsubsection*{关键数据}
\begin{itemize}
  \item 韩国经常账户盈余2026年1-5月达1410亿美元
  \item 2025年韩国企业对美直接投资250亿美元，创历史新高
  \item USD/KRW预计区间1490-1560
  \item 无逆周期因子模型预测USD/CNY中间价6.7936
  \item 逆周期因子影响仅0.6个基点
  \item ILS/KRW自2024年中期以来走强约50\%
  \item 韩国Kospi上涨时外资倾向于减持或增加对冲
\end{itemize}
\begin{figure}[htbp]
\centering
\includegraphics[width=0.88\linewidth]{figures/fig_031.jpg}
\caption{Exhibit 1 - Exhibit 1: Korea's current account surplus has surged with strong exports Exhibit 2: KRW has weakened despite a surge in its terms of trade Exhibit 3: KRW's cheap valuations are comparable to INR and IDR...}
\end{figure}
\begin{figure}[htbp]
\centering
\includegraphics[width=0.88\linewidth]{figures/fig_033.jpg}
\caption{Exhibit 2 - Exhibit 2: KRW has weakened despite a surge in its terms of trade Exhibit 3: KRW's cheap valuations are comparable to INR and IDR... Exhibit 4: ... even though it does not face a similar balance of payments issue What explain}
\end{figure}
\begin{figure}[htbp]
\centering
\includegraphics[width=0.88\linewidth]{figures/fig_059.jpg}
\caption{Figure 2 - Figure 3: Israeli domestic investors drive the shekel through rebalancing and hedging while foreigners have been driving the won}
\end{figure}
\begin{figure}[htbp]
\centering
\includegraphics[width=0.88\linewidth]{figures/fig_061.jpg}
\caption{Figure 2 - Figure 4: ILSKRW is the most expensive cross in the world}
\end{figure}
{\small\color{refgray}本节综合 3 篇研究；涉及 MS、NOM、DB。}
\newpage
\subsection{大宗商品：供给修复与政策博弈重塑中期格局}
\textbf{大宗商品市场正经历从“短缺溢价”向“供给修复+政策反应”的切换：铝因中东复产超预期及印尼产能加速投放，2027年将从短缺转为过剩；铜短期受制于关税不确定性消退与情绪联动，但9月后基本面支撑有望回归；黄金矿企在25\%金价回调中展现出强于历史周期的韧性，但短期成本压力仍存；农产品市场的核心风险已从天气本身转向出口国的预防性保护主义政策。}
\subsubsection*{主流共识}
\begin{itemize}
  \item 全球铝市场2026年仍有小幅缺口，但2027年起因中东产能恢复快于预期及印尼新增产能加速投放，将转为明显过剩。
  \item 铜价短期（7-8月）上涨动能有限，但进入9月后催化剂将重新集结，一年内有望触及15000美元/吨。
  \item 黄金矿企在本轮回调中财务基础显著强于以往周期，创纪录利润率、净现金资产负债表和自律资本配置提供韧性。
  \item 农产品贸易地理集中度极高（前三大出口国控制60-90\%贸易量），政策反应而非天气本身正成为价格的主要驱动因素。
\end{itemize}
\subsubsection*{分歧与条件差异}
\begin{itemize}
  \item 铝市场：MS认为中东产能恢复速度远超预期，2027年过剩超80万吨；而此前市场普遍预期短缺将持续至2027年。
  \item 铜需求：Citi指出5月全球铜终端消费同比下滑10\%几乎完全由去年中国新能源装机基数效应造成，剔除后仍增长约1\%，并非标题所示悲观。
  \item 黄金驱动因素：UBS强调本轮金价驱动力已从实际利率和美元转向央行购金、去美元化及货币信用忧虑，传统宏观指标重新联动但支撑结构未瓦解。
  \item 锂行业：MS认为盛新锂能上半年盈利改善主要依赖锂价高位和印尼产能释放，但上游矿山（木绒锂矿）要到2028年才贡献增量，未来两年处于过渡期。
\end{itemize}
\subsubsection*{机构观点}
\paragraph{UBS} 黄金从2024年四季度到2026年一季度经历约25\%回调，驱动因素已从实际利率和美元转向央行购金、去美元化及货币信用忧虑，中东冲突作为催化剂使传统宏观指标重新联动。金矿股（GDX）较高点回调约35\%，但矿企起点远强于以往周期：创纪录利润率、多数净现金的资产负债表、健康的自由现金流及自律资本配置。短期面临盈利下修、能源成本上升和成本指引不确定性，但全年成本指引大概率维持。纽蒙特自由现金流因资本开支上升环比减半，巴里克燃料成本上升为核心观察点，阿格尼科受柴油成本上升和约6.25亿美元研究流出影响，加拿大马拉蒂克矿坑壁位移事件可能带来每年6-8万盎司产量不确定性。
{\small\color{refgray}数据：金价回调约25\%\par}
{\small\color{refgray}数据：GDX较高点回调约35\%\par}
{\small\color{refgray}数据：矿企多数净现金状态\par}
{\small\color{refgray}数据：自由现金流回报约7\%\par}
{\small\color{refgray}边际变化：与以往周期不同，本轮金价回调中矿企财务韧性显著增强，不再单纯跟随金价下跌，而是具备回购和机会性并购能力。\par}
\paragraph{Citi} 铜价在7-8月难以持续上涨，因市场对美国政府可能征收“232条款”关税的担忧消退，此前存在的COMEX升水面临回落；黄金弱势若持续也会影响铜价情绪。但进入9月后催化剂将重新集结：美联储更可能转向宽松，2027年铜现货市场预计转向约40万吨缺口，AI投资、国防支出和能源转型持续拉动需求。供给约束比需求疲软更值得关注——2026年全球矿山产量预计基本持平，废铜对高价弹性低，中国1-5月废铜进口量同比几乎持平。全球铜终端消费5月同比下滑约10\%，但剔除去年中国新能源装机基数效应后仍增长约1\%。
{\small\color{refgray}数据：一年内铜价目标15000美元/吨\par}
{\small\color{refgray}数据：2027年铜现货市场预计短缺约40万吨\par}
{\small\color{refgray}数据：5月全球铜终端消费剔除基数后增长约1\%\par}
{\small\color{refgray}数据：中国1-5月废铜进口量同比几乎持平\par}
{\small\color{refgray}边际变化：市场此前过度解读5月铜消费同比下滑为需求疲软，Citi指出这几乎完全由基数效应造成，剔除后仍为正增长，供给约束才是中期看多的核心支撑。\par}
\paragraph{GS} 全球农产品贸易地理集中度极高，大豆、玉米、大米、糖和棕榈油前三大出口国控制60-90\%贸易量。价格不确定性的关键来源已从天气本身转向各国出于粮食和能源安全考虑的预防性保护主义政策。当前厄尔尼诺条件已形成，有63\%概率发展为“超级厄尔尼诺”，但2023年印度案例表明，即使最终产量未受显著影响，仅对冲击的担忧就足以触发出口限制并推高价格。2026年上半年能源价格走高可能促使各国提高生物燃料掺混比例，将原本用于出口的糖、棕榈油、玉米转向国内燃料生产，印尼生物柴油掺混率从20\%提至40-50\%即为典型。
{\small\color{refgray}数据：前三大出口国控制60-90\%贸易量\par}
{\small\color{refgray}数据：厄尔尼诺发展为超级厄尔尼诺概率63\%\par}
{\small\color{refgray}数据：印度占全球大米出口约40\%\par}
{\small\color{refgray}数据：印尼生物柴油掺混率从20\%提至40-50\%\par}
{\small\color{refgray}边际变化：市场正从“供给冲击定价”阶段进入“政策反应函数定价”阶段，预防性出口限制和生物燃料政策成为比天气本身更大的价格驱动因素。\par}
\paragraph{MS} 铝市场供需格局正在逆转：中东约350万吨年化产能中断（占全球4-4.8\%）的恢复速度远超预期，大部分产能有望在2027年上半年结束前回归；印尼铝产量2026-2028年分别达86.6万、120万和130万吨，全球其他地区同期新增44万、83.5万和90.5万吨。2027年全球新供应约210万吨（不含中东复产），市场从之前预计的80万吨缺口转为80万吨以上过剩。铝价预计从2026下半年的3150美元/吨降至2027年均值2850美元/吨。盛新锂能上半年净利润预计10-12亿元（去年同期亏损8.41亿元），驱动因素为锂价从6.9万元/吨上涨至13万元/吨及印尼工厂出货量增长，但上游木绒锂矿要到2028年才贡献增量。
{\small\color{refgray}数据：中东产能中断约350万吨/年（全球4-4.8\%）\par}
{\small\color{refgray}数据：2027年全球新供应约210万吨（不含中东复产）\par}
{\small\color{refgray}数据：2027年铝市场从缺口80万吨转为过剩80万吨以上\par}
{\small\color{refgray}数据：铝价2027年均值降至2850美元/吨\par}
{\small\color{refgray}边际变化：市场此前线性外推“短缺”叙事，但MS指出中东产能恢复速度和印尼新增供应节奏同时超出模型假设，2027年铝市场从短缺直接转为过剩。\par}
\subsubsection*{关键数据}
\begin{itemize}
  \item 金价回调约25\%，GDX回调约35\%
  \item 一年内铜价目标15000美元/吨，2027年铜现货预计短缺约40万吨
  \item 前三大出口国控制60-90\%农产品贸易量，厄尔尼诺发展为超级厄尔尼诺概率63\%
  \item 中东铝产能中断约350万吨/年（全球4-4.8\%），2027年铝市场从缺口80万吨转为过剩80万吨以上
  \item 铝价2027年均值降至2850美元/吨
  \item 盛新锂能1H26净利润10-12亿元（去年同期亏损8.41亿元），锂价从6.9万元/吨上涨至13万元/吨
  \item 印尼铝产量2026年86.6万吨、2027年120万吨、2028年130万吨
  \item 5月全球铜终端消费剔除基数后仍增长约1\%
\end{itemize}
\begin{figure}[htbp]
\centering
\includegraphics[width=0.88\linewidth]{figures/fig_051.jpg}
\caption{Figure 1 - \# Copper upside from September, medium-term bullish view intact We expect renewed upside for copper prices after the summer as we look for prices to average \textbackslash{}\$14,500/t in 4Q'26 and to touch \textbackslash{}\$15k/t within a year. We think copper will struggle to rally sustaina}
\end{figure}
\begin{figure}[htbp]
\centering
\includegraphics[width=0.88\linewidth]{figures/fig_066.jpg}
\caption{Exhibit 1 - Exhibit 2: Consistent With Our Commodity Control Cycle Framework, Concentrated Supply Increases Disruption Risk}
\end{figure}
\begin{figure}[htbp]
\centering
\includegraphics[width=0.88\linewidth]{figures/fig_122.jpg}
\caption{Exhibit 4 - Exhibit 5: The aluminum forward curve in backwardation discouraged inventory holding Exhibit 6: US imports fell sharply in April Exhibit 7: The YoY change in annualised Middle East production has been smaller than expected vers}
\end{figure}
\begin{figure}[htbp]
\centering
\includegraphics[width=0.88\linewidth]{figures/fig_069.jpg}
\caption{Exhibit 2 - Exhibit 4: Ahead of 2023 El Nino Concerns, India Imposed a Precautionary Rice Export Ban, Despite Domestic Production Ultimately Proving Resilient Thailand white rice \$5\textbackslash{}\%\$ is the benchmark price for Asian rice prices.}
\end{figure}
{\small\color{refgray}本节综合 5 篇研究；涉及 UBS、Citi、GS、MS。}
\newpage
\subsection{AI基础设施与模型竞争：从算力军备到效率优先}
\textbf{AI投资正从‘算力军备’转向‘效率优先’，基础设施层（记忆体长协、HDD、硅片）因AI需求出现结构性量价齐升，模型层中国开源模型逼近‘足够好’临界点，成本效率成为竞争核心。}
\subsubsection*{主流共识}
\begin{itemize}
  \item AI需求正从超大规模云厂商向主权国家、企业本地部署扩散，非超大规模客户市场最终更大。
  \item 记忆体长协因财务前置和AI买家信用，供应商获得前所未有的节奏保护，违约成本被前置锁定。
  \item 近线HDD供需紧张至少持续至2027年，收入增速（+37\%）远超出货量增速（+12\%），ASP结构性提升。
  \item 中国AI开源模型在智能表现上已逼近全球闭源模型，成本-效果比驱动企业采用加速。
\end{itemize}
\subsubsection*{分歧与条件差异}
\begin{itemize}
  \item 物理AI商业化已发生，但规模化瓶颈（数据稀缺、人才、电池、部署成本）是运营问题而非技术问题，短期回报来自专用AMR而非人形机器人。
  \item Netflix下半年用户增长依赖广告超预期而非内容爆发，世界杯冲击打乱季节性节奏，人均观看时长持续承压。
  \item 小米AI价值（保守320亿元）未被市场定价，市场焦点过度集中于电动车出货量和芯片涨价。
\end{itemize}
\subsubsection*{机构观点}
\paragraph{Bernstein} 本轮记忆体长协与历史案例存在结构性差异：财务担保采用‘后端加重’设计，合同后期覆盖率升至75-100\%，违约成本被前置锁定；AI买家信用等级高、需求结构性而非周期性，供应商获得前所未有的节奏保护。CoreWeave二季度为桥接季，真正的业绩拐点在下半年电力上线后，但执行容错空间极小，利润率、利息负担是关键。Netflix上半年用户净增仅300-400万，世界杯打乱季节性节奏，下半年内容缺乏爆款，收入增长更依赖广告超预期。
{\small\color{refgray}数据：记忆体长协后期财务担保覆盖率75-100\%\par}
{\small\color{refgray}数据：CoreWeave 2Q26 AOI指引3000-9000万美元，下半年需7.9-10.5亿\par}
{\small\color{refgray}数据：Netflix 1H26总观看时长低于960亿小时\par}
{\small\color{refgray}数据：Netflix 1H26净增用户约300-400万\par}
{\small\color{refgray}边际变化：记忆体长协从‘法律承诺’转向‘真实抵押资产’，违约成本被前置锁定；CoreWeave从合同驱动转向执行验证；Netflix从内容驱动转向广告驱动。\par}
\paragraph{Citi} 英伟达最被低估的护城河是路线图的确定性和能效提升：Blackwell实现比Hopper 25倍能效提升，每GW算力成本从30-40亿美元向100亿美元演进，定价权从‘算力性能’向‘算力效率’迁移。物理AI商业化已发生，人形机器人定价已低于人类劳动力成本，投资回收期约1年，但规模化瓶颈在于数据稀缺（需要‘专家token’）、人才约束、电池续航和高部署成本，短期回报来自专用AMR。
{\small\color{refgray}数据：Blackwell比Hopper能效提升25倍\par}
{\small\color{refgray}数据：每GW算力成本从30-40亿美元向100亿美元演进\par}
{\small\color{refgray}数据：人形机器人投资回收期约1年\par}
{\small\color{refgray}数据：某物理AI公司ARR接近3亿美元，连续3年营收翻3倍\par}
{\small\color{refgray}边际变化：英伟达从‘卖芯片’转向‘卖能效和系统’；物理AI从概念验证进入商业部署，但规模化瓶颈从技术问题转为运营问题。\par}
\paragraph{MS} 市场对小米的定价中几乎未反映AI业务潜在价值，保守估值已达320亿元，而芯片涨价对智能手机的影响仅约210亿元。电动车出货量预测下调（2026年从58万至50万辆），但长期结构性机会被低估。市场焦点过度集中于短期周期性问题，忽略了小米从硬件制造商向AI生态企业转型的价值。
{\small\color{refgray}数据：小米AI业务保守估值320亿元\par}
{\small\color{refgray}数据：芯片涨价对智能手机估值影响约210亿元\par}
{\small\color{refgray}数据：2026年电动车出货量预测从58万下调至50万辆\par}
{\small\color{refgray}数据：熊市情景概率从10\%上调至20\%\par}
{\small\color{refgray}边际变化：小米估值焦点从电动车出货量转向AI价值重估，市场尚未计入AI生态转型溢价。\par}
\paragraph{NOM} 天孚通信1.6T光引擎拐点或在2026下半年，大型AI客户帮助获取更多200G EML芯片供应，1.6T光引擎业务2026-2028年收入CAGR达48\%，到2028财年贡献总销售额55\%。市场对玻璃基板等新技术冲击的担忧过度，商业化时间表尚远。鼎龙半导体材料业务进入规模化阶段，CMP抛光液/垫、光刻胶三大产品线同步放量，1H26净利润同比增64-74\%，扣非后增长超80\%。
{\small\color{refgray}数据：天孚1.6T光引擎2026-2028年收入CAGR 48\%\par}
{\small\color{refgray}数据：天孚1.6T光引擎2028财年贡献总销售额55\%\par}
{\small\color{refgray}数据：鼎龙1H26净利润同比增64-74\%，扣非后增长超80\%\par}
{\small\color{refgray}数据：鼎龙CMP抛光液6月单月销售突破4000万元\par}
{\small\color{refgray}边际变化：天孚从200G EML短缺担忧转向供给缓解拐点；鼎龙从‘验证通过’走向‘规模盈利’，材料环节国产化进入财务报表可验证阶段。\par}
\paragraph{GS} 近线HDD供需紧张至少持续至2027年，2026年行业收入预计从223亿跃升至305亿美元（+37\%），ASP从180美元跳升至236美元（+30\%），涨价从‘缺货结果’转为主动产能策略。硅片长约从5-7年延长至10年，存储客户为AI需求主动锁仓，美光向环球晶圆提供5亿美元战略融资，下一轮长约议价权向卖方倾斜。中国AI开源模型在智能表现上逼近全球闭源模型，预计2026年底行业ARR达100亿美元，2030年达1250亿美元，MiniMax以0.4T参数达到DeepSeek V4 Pro 1.6T参数类似效果，成本效率优势突出。
{\small\color{refgray}数据：2026年近线HDD出货量7560万台，同比+12\%\par}
{\small\color{refgray}数据：2026年HDD行业收入305亿美元，ASP 236美元\par}
{\small\color{refgray}数据：美光向环球晶圆提供5亿美元战略融资，签署10年长约\par}
{\small\color{refgray}数据：中国AI模型行业2026年底ARR预计100亿美元，2030年1250亿美元\par}
{\small\color{refgray}边际变化：HDD从‘量价齐跌’转向‘量价齐升’，涨价成为主动策略；硅片长约从5-7年延长至10年，存储客户从需求波动放大器变为稳定性背书者；中国AI模型从‘成本效率突破’转向‘智能水平逼近’，开源模型启动正向数据飞轮。\par}
\subsubsection*{关键数据}
\begin{itemize}
  \item 记忆体长协后期财务担保覆盖率75-100\%
  \item Blackwell比Hopper能效提升25倍
  \item 2026年HDD行业收入305亿美元，ASP 236美元
  \item 美光向环球晶圆提供5亿美元战略融资，签署10年长约
  \item 中国AI模型行业2026年底ARR预计100亿美元
  \item MiniMax M3模型0.4T参数达到DeepSeek V4 Pro 1.6T参数类似效果
  \item 小米AI业务保守估值320亿元
  \item 天孚1.6T光引擎2026-2028年收入CAGR 48\%
  \item 鼎龙1H26净利润同比增64-74\%，扣非后增长超80\%
  \item Netflix 1H26净增用户约300-400万
\end{itemize}
\begin{figure}[htbp]
\centering
\includegraphics[width=0.88\linewidth]{figures/fig_001.jpg}
\caption{EXHIBIT 1 - EXHIBIT 4: CRWV 1Q26 Post-Earnings Stock Reaction}
\end{figure}
\begin{figure}[htbp]
\centering
\includegraphics[width=0.88\linewidth]{figures/fig_021.jpg}
\caption{Exhibit 3 - EXHIBIT 1: Engagement is expected to be under 96B hours in 1H26, lower than 2H25, which had a strong content slate Half-year total is expected to be under 96B EXHIBIT 2: Netflix experienced declining engagement in Q2, and did n}
\end{figure}
\begin{figure}[htbp]
\centering
\includegraphics[width=0.88\linewidth]{figures/fig_041.jpg}
\caption{Exhibit 1 - Exhibit 1: Our base/bull/bear cases for Xiaomi's Mi EV product roadmap and sales volume forecasts 2024 total volume \& mkt share\% 2025A total volume \& mkt share\% 2026E total volume \& mkt share\% 2027E total volume \& mkt share\%}
\end{figure}
\begin{figure}[htbp]
\centering
\includegraphics[width=0.88\linewidth]{figures/fig_071.jpg}
\caption{Exhibit 1 - Exhibit 1: Unit economics for China's value-for-money agentic model and top performing coding model, vs. hypothetical economics for global SOTA model UNIT ECONOMICS OF TWO-TIERED CHINA AI MODELS (PER 1M BLENDED TOKENS) \#\# IN ABSOLU}
\end{figure}
{\small\color{refgray}本节综合 12 篇研究；涉及 Bernstein、Citi、MS、NOM、GS。}
\newpage
\subsection{AI供应链：从PCB到服务器机柜的产能爬坡与地理重构}
\textbf{AI硬件供应链正经历双重结构性变化：一是PCB厂商在东南亚布局高端产能以服务海外客户，实现产品组合升级与利润修复；二是AI服务器机柜出货量保持环比增长，全年预测不变，下半年仍有加速预期。}
\subsubsection*{主流共识}
\begin{itemize}
  \item AI服务器和高速交换机推动PCB规格升级，单机价值量提升，全球AI PCB需求处于上升通道。
  \item PCB厂商在泰国等东南亚地区扩产，主要服务海外客户的AI相关高端需求，并非简单的成本套利。
  \item AI服务器机柜出货量2026年全年预测70-80K台不变，下半年环比增速将加快。
  \item PCB钻针等上游耗材的订单增长可作为AI PCB产业链景气的先行指标。
\end{itemize}
\subsubsection*{分歧与条件差异}
\begin{itemize}
  \item PCB新产能爬坡初期对利润率的影响：GS认为随着规模扩Daiwa良率提升，盈利将逐步改善；但短期折旧和材料成本可能压制利润。
  \item ODM出货节奏分化：鸿海二季度略超预期，而广达和纬创部分出货递延至下半年，市场对短期节奏与全年目标的关系存在不同解读。
\end{itemize}
\subsubsection*{机构观点}
\paragraph{GS} GS在泰国科技之旅中走访多家PCB及钻针厂商，确认AI PCB需求正从预期走向可追踪的订单与扩产。PCB厂商在泰国追加投资（如某厂追加8000万美元、另一厂33亿元计划），专攻海外客户AI服务器和光模块高端PCB，产品定价和利润率显著高于常规PCB。同时，PCB钻针供应商DTech泰国月产能已达600万支，扩产节奏与客户海外工厂同步，验证需求真实落地。GS认为，AI PCB规格升级、单机价值量提升、全球需求扩张三重逻辑同步兑现，中国头部PCB和CCL供应商将受益。
{\small\color{refgray}数据：某PCB厂商泰国工厂2026上半年净利润预计同比增长85\%-95\%\par}
{\small\color{refgray}数据：某PCB厂商追加8000万美元投入泰国高端PCB产能\par}
{\small\color{refgray}数据：DTech泰国工厂2026年4月月产能达600万支，计划增资不超过3亿元\par}
{\small\color{refgray}数据：全球AI PCB市场规模2027年预计达183亿美元，2025-2027年CAGR 140\%\par}
{\small\color{refgray}边际变化：从关注国内PCB产能转向东南亚高端产能布局，强调海外产能对锁定AI订单的战略价值，以及钻针订单作为产业链景气先行指标的作用。\par}
\paragraph{MS} MS 6月AI服务器机柜出货追踪显示，GB200/300 NVL72机柜单月产出约8000台，环比增长5\%，全年70-80K台预测不变。鸿海6月出货3300台，二季度累计10300台，略超预期，三季度指引维持增长；广达出货2300-2400台，纬创1200-1300台，两者二季度略低于预期，但部分出货仅递延至下半年，并非订单消失。MS认为，下半年月均出货需达8-9K台才能完成全年目标，环比增速需进一步提升，但整体增长趋势未变。
{\small\color{refgray}数据：6月AI服务器机柜出货约8000台，环比+5\%\par}
{\small\color{refgray}数据：鸿海6月出货3300台，二季度累计10300台，环比+21\%\par}
{\small\color{refgray}数据：广达6月出货2300-2400台，二季度累计6200-6300台\par}
{\small\color{refgray}数据：纬创6月出货1200-1300台，二季度累计3900-4000台\par}
{\small\color{refgray}边际变化：从关注单月出货绝对量转向强调出货时间平移而非订单消失，以及下半年环比加速预期。\par}
\subsubsection*{关键数据}
\begin{itemize}
  \item 某PCB厂商泰国工厂2026上半年净利润预计同比增长85\%-95\%
  \item 某PCB厂商追加8000万美元投入泰国高端PCB产能
  \item DTech泰国工厂2026年4月月产能达600万支，计划增资不超过3亿元
  \item 全球AI PCB市场规模2027年预计达183亿美元，2025-2027年CAGR 140\%
  \item 6月AI服务器机柜出货约8000台，环比+5\%
  \item 鸿海6月出货3300台，二季度累计10300台，环比+21\%
  \item 全年AI服务器机柜出货预测70-80K台，为2025年约29K台的2倍以上
\end{itemize}
\begin{figure}[htbp]
\centering
\includegraphics[width=0.88\linewidth]{figures/fig_114.jpg}
\caption{Exhibit 1 - Exhibit 1: Industry-wide GB200/300 NVL72-equivalent monthly rack output, by major ODMs (2025) Exhibit 2: Industry-wide GB200/300 NVL72-equivalent monthly rack output, by major ODMs (2026) GB200/300 NVL72 racks by major ODMs (000s)}
\end{figure}
\begin{figure}[htbp]
\centering
\includegraphics[width=0.88\linewidth]{figures/fig_115.jpg}
\caption{Exhibit 2 - Exhibit 2: Industry-wide GB200/300 NVL72-equivalent monthly rack output, by major ODMs (2026) GB200/300 NVL72 racks by major ODMs (000s) GB200/300 NVL72 racks by major ODMs (000s) \#\# Exhibit 3: GB200/300 NVL72 racks by major OD}
\end{figure}
\begin{figure}[htbp]
\centering
\includegraphics[width=0.88\linewidth]{figures/fig_116.jpg}
\caption{Exhibit 3 - Exhibit 3: GB200/300 NVL72 racks by major ODMs (000s) Exhibit 4: GB200/300 rack supply share, by major ODMs (2025) GB200/300 NVL72-equivalent rack supply share (2025)}
\end{figure}
\begin{figure}[htbp]
\centering
\includegraphics[width=0.88\linewidth]{figures/fig_117.jpg}
\caption{Exhibit 4 - Exhibit 4: GB200/300 rack supply share, by major ODMs (2025) GB200/300 NVL72-equivalent rack supply share (2025) Exhibit 5: GB200/300 rack supply share, by major ODMs (2026) GB200/300 NVL72-equivalent rack supply share (2026) \#}
\end{figure}
{\small\color{refgray}本节综合 4 篇研究；涉及 GS、MS。}
\newpage
\subsection{工业与机器人：人形机器人量产加速，电力设备订单兑现}
\textbf{人形机器人正从实验室走向工厂，量产节奏超预期，带动中国供应链企业（恒立液压、绿的谐波）收入占比快速提升；同时，数据中心电力需求爆发，特锐德订单已超去年全年，而韩国核电EPC订单时间线因政治因素推迟1-2年。}
\subsubsection*{主流共识}
\begin{itemize}
  \item 人形机器人量产节奏快于预期，2026年Q3-Q4产量将从每周约100台爬坡至约1.9万台，2027年目标20万台。
  \item 中国零部件供应商（恒立液压、绿的谐波）在人形机器人供应链中角色日益重要，收入占比将从2026年的约1\%提升至2027年的约8\%。
  \item 数据中心电力基础设施需求强劲，高压预装式变电站订单增长显著，特锐德年初至今订单已超2025年全年。
\end{itemize}
\subsubsection*{分歧与条件差异}
\begin{itemize}
  \item 人形机器人量产节奏：Citi基于供应链调研认为2026年Q3可达每周1000台，但部分市场参与者预期更为保守。
  \item 韩国核电EPC订单时间：JPM将美国项目订单时间从2028年推迟至2029年，但市场对此已有预期，分歧在于推迟幅度是否已被充分定价。
  \item 恒立液压短期汇兑损失：Citi认为约2亿元汇兑损失是短期噪音，不反映经营基本面，但部分投资者担忧汇率波动对净利润的持续影响。
\end{itemize}
\subsubsection*{机构观点}
\paragraph{Citi} 恒立液压人形机器人业务放量节奏超预期，美国客户产量将从当前每周约100台爬坡至2026年Q3约1000台、Q4约1.9万台，2027年目标20万台。预计人形机器人收入占比从2026年约1\%升至2027年约8\%。核心业务挖掘机部件需求强劲，2Q26营收同比增长约30\%，毛利率维持44\%，但约2亿元汇兑损失拖累净利润增速仅9\%。川崎重工调整中国液压件布局，恒立市场份额持续提升。
{\small\color{refgray}数据：2026年Q3/Q4美国客户人形机器人产量预期约1000台/1.9万台每周\par}
{\small\color{refgray}数据：2027年产量目标20万台\par}
{\small\color{refgray}数据：人形机器人收入占比从2026年约1\%升至2027年约8\%\par}
{\small\color{refgray}数据：2Q26营收同比增长约30\%，毛利率约44\%\par}
{\small\color{refgray}边际变化：人形机器人量产节奏从千台级别跃升至20万台级别，供应链角色从验证阶段进入实质放量阶段。\par}
\paragraph{Citi} 绿的谐波正从工业机器人零部件供应商转型为人形机器人精密组件核心玩家，通过两个合资动作（与敏实在美国成立人形机器人合资公司、与SKF在中国成立精密部件合资公司）锁定供应链地位。与SKF合资生产交叉滚子轴承、柔性轴承等，一台人形机器人需60-70个轴承，占BOM成本5\%-15\%。潜在美国客户量产在即，预计2026年Q3约1000台、Q4约1.9万台，2027年增长10倍至约20万台。
{\small\color{refgray}数据：一台人形机器人需60-70个轴承单元\par}
{\small\color{refgray}数据：轴承成本占BOM 5\%-15\%\par}
{\small\color{refgray}数据：2026年Q3/Q4美国客户产量预期约1000台/1.9万台每周\par}
{\small\color{refgray}数据：2027年产量目标20万台\par}
{\small\color{refgray}边际变化：绿的谐波从工业机器人谐波减速器供应商转型为人形机器人关节精密组件供应商，合资公司锁定客户与产能。\par}
\paragraph{JPM} 韩国核电EPC行业2026年上半年新订单几乎为零，但市场已有预期。海外核电项目时间线整体推迟1-2年：美国项目从2028年推迟至2029年（受韩美2000亿美元基建投资分配节奏影响），保加利亚项目推迟至2027年。下调斗山能源、KEPCO E\&C、现代建设2026-28年订单预估3\%-31\%。非核电业务（斗山美国蒸汽轮机、现代建设国内重建项目）填补空白期，2Q业绩预计符合预期。
{\small\color{refgray}数据：美国核电项目订单时间从2028年推迟至2029年\par}
{\small\color{refgray}数据：保加利亚项目推迟至2027年\par}
{\small\color{refgray}数据：斗山能源订单预估下调5\%-17\%\par}
{\small\color{refgray}数据：KEPCO E\&C订单预估下调3\%-31\%\par}
{\small\color{refgray}边际变化：核电订单时间线因政治因素推迟1-2年，市场关注点从“谁能拿到订单”转向“订单何时转化为收入”。\par}
\paragraph{JPM} 特锐德数据中心预装式变电站订单显著增长，年初至今新订单4-5亿元已超2025年全年，管理层指引全年超8亿元。在大型数据中心（50MW以上）110kV及以上高压预装式变电站领域市占率约50\%。与伊顿合作整合固态变压器技术，形成110/220kV转800V直流的一站式模块化方案。中东订单有望翻倍以上增长，沙特能源公司预计Q3招标40-60亿元。2025-2028年盈利复合增速超20\%。
{\small\color{refgray}数据：年初至今数据中心新订单4-5亿元，超2025年全年\par}
{\small\color{refgray}数据：管理层指引全年数据中心订单超8亿元\par}
{\small\color{refgray}数据：大型数据中心高压预装式变电站市占率约50\%\par}
{\small\color{refgray}数据：已为腾讯两个100MW数据中心供货\par}
{\small\color{refgray}边际变化：数据中心订单从预期转向兑现，年初至今已超去年全年，中东市场弹性可期。\par}
\subsubsection*{关键数据}
\begin{itemize}
  \item 人形机器人美国客户2026年Q3/Q4产量预期约1000台/1.9万台每周，2027年目标20万台
  \item 恒立液压人形机器人收入占比从2026年约1\%升至2027年约8\%
  \item 绿的谐波人形机器人收入占比预计从2025年约20\%升至2026年40\%-50\%
  \item 特锐德年初至今数据中心新订单4-5亿元，超2025年全年
  \item 特锐德大型数据中心高压预装式变电站市占率约50\%
  \item 韩国核电美国项目订单时间从2028年推迟至2029年
  \item 斗山能源订单预估下调5\%-17\%
  \item 恒立液压2Q26汇兑损失约2亿元
\end{itemize}
\begin{figure}[htbp]
\centering
\includegraphics[width=0.88\linewidth]{figures/fig_029.jpg}
\caption{Figure 5 - Figure 5. Hengli: US humanoid robot client's production estimate © 2026 Citi Inc. No redistribution without Citi's written permission. Figure 6. Hengli: Humanoid robot revenue and as \% of total revenue © 2026 Citi Inc. No redistribution without Citi's written }
\end{figure}
\begin{figure}[htbp]
\centering
\includegraphics[width=0.88\linewidth]{figures/fig_027.jpg}
\caption{Figure 3 - Figure 3. USDCNY and its changes © 2026 Citi Inc. No redistribution without Citi's written permission. Figure 4. Hengli: Financial costs vs. USDCNY © 2026 Citi Inc. No redistribution without Citi's written permission. Figure 5. Hengli: US humanoid robot client}
\end{figure}
\begin{figure}[htbp]
\centering
\includegraphics[width=0.88\linewidth]{figures/fig_106.jpg}
\caption{Figure 1 - Figure 1: Doosan: Share price vs orderbook Figure 2: Doosan: Price/orderbook trend \# Investment Thesis, Valuation and Risks}
\end{figure}
\begin{figure}[htbp]
\centering
\includegraphics[width=0.88\linewidth]{figures/fig_108.jpg}
\caption{Figure 4 - Figure 4: KEPCO E\&C: Price/orderbook trend Figure 3: KEPCO E\&C: Share price vs orderbook Won, Wbn (RHS) \# Investment Thesis, Valuation and Risks KEPCO E\&C (Overweight; Price Target: W180,000)}
\end{figure}
{\small\color{refgray}本节综合 4 篇研究；涉及 Citi、JPM。}
\newpage
\subsection{消费与零售：结构分化加剧，品牌替代与区域集中成为新主线}
\textbf{全球消费市场正经历从总量增长向结构分化的转变：日本消费增长高度集中于东京-福冈黄金走廊的6个核心县；中国PC/平板市场总量承压但品牌替代效应显著；泡泡玛特IP热度见顶，增长模式面临调整；而极兔在拉美等低渗透市场的扩张则验证了电商物流的蓝海逻辑。}
\subsubsection*{主流共识}
\begin{itemize}
  \item 日本消费增长动力集中在东京-福冈黄金走廊的6个核心县，其余地区贡献有限。
  \item 中国PC与平板市场总量持续下滑，但品牌间替代效应显著，高端品牌逆势增长。
  \item 泡泡玛特核心IP Labubu 4.0二级市场折价表明IP热度见顶，增长模式从稀缺转向调整。
  \item 极兔已从建设期进入规模与效率回报期，海外市场占比提升至32.3\%。
\end{itemize}
\subsubsection*{分歧与条件差异}
\begin{itemize}
  \item 日本消费增长：Bernstein强调地理集中度，认为聚焦6县可获更快增长；而市场普遍仍以全国视角布局。
  \item 中国PC市场：MS认为总量下滑但苹果等品牌通过替代效应实现增长；部分观点仍等待总量拐点。
  \item 泡泡玛特：DB认为IP周期已过，下半年营收可能同比下滑；市场此前定价隐含持续增长预期。
  \item 乳制品行业：MS认为原奶价格稳定后龙头优势显现；但定价能力恢复仍需时间，利润爬坡缓慢。
\end{itemize}
\subsubsection*{机构观点}
\paragraph{Bernstein} 日本消费增长高度集中于东京-福冈黄金走廊的6个核心县（东京、千叶、神奈川、爱知、大阪、福冈），这些县聚集了44\%的劳动年龄人口（约3000万人），家庭数量增速（2019-2024年CAGR 1.3\%）是其余41个县的两倍多。收入增长同样分化：核心县收入CAGR 2.8\%，其余地区仅1.3\%。千叶消费支出CAGR 5.1\%，远超东京的3.2\%，其中非酒精饮料消费增速10.4\%，是6县平均的4倍。聚焦这6个县可覆盖近一半购买力，获得比全国快约30\%的消费支出增长。
{\small\color{refgray}数据：6个核心县占日本44\%劳动年龄人口（约3000万人）\par}
{\small\color{refgray}数据：核心县家庭数量CAGR 1.3\%，其余地区0.6\%\par}
{\small\color{refgray}数据：核心县收入CAGR 2.8\%，其余地区1.3\%\par}
{\small\color{refgray}数据：千叶消费支出CAGR 5.1\%，东京3.2\%\par}
{\small\color{refgray}边际变化：从全国均匀布局转向聚焦黄金走廊6县，量化了地理集中度对消费增长的差异化影响。\par}
\paragraph{Citi} 国巨2Q26营收445亿新台币，创历史新高，较市场预期高6\%，增长来自标准品和特殊品同步走强，而非单一产品线驱动。BB值持续高于1.0，AI相关订单仍在加速，上半年营收826亿新台币已占市场全年预测的47\%。随着下半年新一代AI GPU和ASIC平台放量、产品定价调整加速，全年数字存在明显上修空间。需求广度是这轮周期的真正变量，下游工业、车用、消费电子等多个终端同时回温。
{\small\color{refgray}数据：2Q26营收445亿新台币，同比+36\%，环比+16\%\par}
{\small\color{refgray}数据：6月营收154亿新台币，同比+39\%\par}
{\small\color{refgray}数据：上半年营收826亿新台币，占市场全年预测47\%\par}
{\small\color{refgray}数据：BB值持续高于1.0\par}
{\small\color{refgray}边际变化：从过去依赖MLCC单一产品线驱动，转向标准品和特殊品同步增长，需求基础更广泛。\par}
\paragraph{DB} 泡泡玛特2Q26集团营收预计同比仅增1.6\%，较1Q26的75-80\%增速大幅回落。核心IP Labubu 4.0上市一周后二级市场标准款跌至约50元人民币，远低于159元零售价，折价幅度超过其他IP，表明藏家和黄牛热情降温。海外市场同步调整，收入环比降20\%，同比降19\%，印尼、新加坡、泰国做福袋清库存，欧美第三方渠道打8折。618期间天猫店满减活动，Labubu世界杯系列被塞进福袋当赠品，供需关系已从抢不到转为打折卖。
{\small\color{refgray}数据：2Q26集团营收预计同比+1.6\%，1Q26为75-80\%\par}
{\small\color{refgray}数据：Labubu 4.0二级市场标准款跌至约50元，零售价159元\par}
{\small\color{refgray}数据：海外收入环比降20\%，同比降19\%\par}
{\small\color{refgray}数据：618期间Labubu世界杯系列被用作福袋赠品\par}
{\small\color{refgray}边际变化：从供不应求的稀缺阶段转向需求节奏变化，IP热度见顶，海外扩张同步放缓。\par}
\paragraph{JPM} 极兔2Q26集团包裹量91.8亿件，同比+24.2\%，日均首次突破1亿件，标志公司从建设期进入规模与效率回报期。海外包裹占比32.3\%，其中拉美市场包裹量2.1亿件，同比+137\%，连续5个季度加速。东南亚包裹27.5亿件，同比+63.2\%，竞争从价格转向成本、覆盖和可靠性。中国包裹62.1亿件，同比+10.6\%，跑赢行业5\%增速，靠网络升级而非价格竞争。20亿港元回购计划，按当前股价约2.4\%回购收益率。
{\small\color{refgray}数据：2Q26集团包裹量91.8亿件，日均首次突破1亿件\par}
{\small\color{refgray}数据：海外包裹占比32.3\%\par}
{\small\color{refgray}数据：拉美包裹量2.1亿件，同比+137\%\par}
{\small\color{refgray}数据：东南亚包裹27.5亿件，同比+63.2\%\par}
{\small\color{refgray}边际变化：从建设+亏损阶段迈入规模、效率与资本回报新周期，海外市场成为增长新引擎。\par}
\paragraph{MS} 蒙牛2026年前5个月收入增速接近高单位数，鲜奶预期增长25-30\%，奶粉超30\%，UHT已恢复正增长但由量驱动而非提价。原奶价格环比稳定，行业竞争趋于理性，但全年经营利润率预计同比持平，利润爬坡需3-5年。高端产品特仑苏表现优于基础UHT，产品结构改善带来利润弹性。鲜奶产品可能面临税收调整，但占蒙牛总收入仅中单位数，影响有限。
{\small\color{refgray}数据：前5个月收入增速接近高单位数（7-9\%）\par}
{\small\color{refgray}数据：鲜奶预期上半年增长25-30\%\par}
{\small\color{refgray}数据：奶粉上半年增速超30\%\par}
{\small\color{refgray}数据：UHT恢复正增长但定价偏弱\par}
{\small\color{refgray}边际变化：从原奶价格波动和行业亏损阶段转向成本稳定、竞争理性，但利润修复缓慢。\par}
\paragraph{MS} 中国5月PC出货量同比下滑14\%至289万台，平板下滑7\%至281万台，但品牌替代效应显著。苹果PC同比+45\%，年初至今累计+48\%，从华为（-39\%）和戴尔（-39\%）手中抢走份额。平板市场联想同比+42\%，华为+7\%，苹果iPad同比-13\%。总量调整下，高端替代窗口已打开，企业级和高端消费群体加速向特定品牌集中。
{\small\color{refgray}数据：5月PC出货289万台，同比-14\%\par}
{\small\color{refgray}数据：苹果PC同比+45\%，年初至今累计+48\%\par}
{\small\color{refgray}数据：华为PC同比-39\%，戴尔同比-39\%\par}
{\small\color{refgray}数据：平板出货281万台，同比-7\%\par}
{\small\color{refgray}边际变化：从等待总量拐点转向关注品牌替代效应，高端品牌在总量下滑中逆势增长。\par}
\subsubsection*{关键数据}
\begin{itemize}
  \item 日本6个核心县占44\%劳动年龄人口，消费增速比全国快30\%
  \item 国巨2Q26营收445亿新台币创历史新高，BB值持续>1.0
  \item Labubu 4.0二级市场跌至50元，远低于159元零售价
  \item 极兔日均包裹量首破1亿件，拉美包裹量同比+137\%
  \item 蒙牛鲜奶预期增长25-30\%，全年经营利润率持平
  \item 中国5月PC同比-14\%，苹果PC逆势+45\%
  \item 泡泡玛特2Q26营收增速从75-80\%骤降至1.6\%
  \item 极兔海外包裹占比32.3\%，20亿港元回购计划
\end{itemize}
\begin{figure}[htbp]
\centering
\includegraphics[width=0.88\linewidth]{figures/fig_016.jpg}
\caption{EXHIBIT 1 - EXHIBIT 1: \$70\textbackslash{}\%\$ of Japan's population lives in a narrow band stretching from Tokyo to Fukuoka As of Oct 2024 Around \$70\textbackslash{}\%\$ of Japan's population, totaling c. 87mn people, live in a narrow band stretching from Tokyo in the East}
\end{figure}
\begin{figure}[htbp]
\centering
\includegraphics[width=0.88\linewidth]{figures/fig_056.jpg}
\caption{Figure 1 - Laura (Chia Yi) Chen Figure 1. Yageo – Quarterly sales trend © 2026 Citi Inc. No redistribution without Citi's written permission. Figure 2. Yageo – Monthly sales trend © 2026 Citi Inc. No redistribution without Citi's written permission.}
\end{figure}
\begin{figure}[htbp]
\centering
\includegraphics[width=0.88\linewidth]{figures/fig_058.jpg}
\caption{Figure 2 - Figure 2: Pop Mart: EBIT Margin Breakdown by region}
\end{figure}
\begin{figure}[htbp]
\centering
\includegraphics[width=0.88\linewidth]{figures/fig_102.jpg}
\caption{Figure 1 - Figure 1: J\&T Parcel Volume Group Figure 2: J\&T Parcel Volume – Southeast Asia Figure 3: J\&T Parcel Volume – New Market}
\end{figure}
{\small\color{refgray}本节综合 6 篇研究；涉及 Bernstein、Citi、DB、JPM、MS。}
\newpage
\subsection{医疗与生物科技：从短期业绩到长期管线，分化与重估并存}
\textbf{大型制药公司正进入‘业绩底牌清晰、但重估叙事模糊’的阶段，短期成本控制带来的超预期难以驱动估值修复，真正的重估动力来自2027年后的管线进展和资本配置战略；与此同时，中国生物科技公司通过差异化研发和BD交易开始获得全球认可，而脑机接口行业正从技术竞赛转向临床疗效比拼，商业化节奏慢于预期。}
\subsubsection*{主流共识}
\begin{itemize}
  \item 大型制药公司短期业绩稳健，但超预期更多来自成本控制而非销售增长，对股价推动作用有限。
  \item 中国生物科技公司通过对外授权交易验证了研发平台价值，差异化竞争（如剂型创新）成为关键。
  \item 脑机接口行业进入临床验证阶段，技术先进性不等于商业成功，临床疗效和商业化路径才是分水岭。
\end{itemize}
\subsubsection*{分歧与条件差异}
\begin{itemize}
  \item 赛诺菲的估值修复时机：MS认为重估动力集中在2027年，当前折价不会很快收窄；而部分市场观点可能期待更早的催化剂。
  \item 中国生物制药的估值逻辑：GS认为其正从‘仿制药企业’向‘创新药平台’切换；但市场仍对其集采不确定性和管线风险存疑。
  \item 脑机接口的商业化路径：JPM认为非侵入式和半侵入式可能占据最大市场；而部分观点仍聚焦于侵入式技术的高信号精度优势。
\end{itemize}
\subsubsection*{机构观点}
\paragraph{MS} 赛诺菲Q2业绩稳健，但超预期主要来自研发费用低于预期（约5\%），推动营业利润比共识高3\%，业务每股收益1.95欧元（比共识高4\%）。全年EPS指引大概率上调至低双位数增长，但市场已部分消化。真正的估值重估动力——新药管线进展和资本配置战略——集中在2027年，当前折价不会很快收窄。
{\small\color{refgray}数据：Q2集团销售额约108亿欧元，与市场预期持平\par}
{\small\color{refgray}数据：度普利尤单抗预计44.3亿欧元，符合预期\par}
{\small\color{refgray}数据：研发费用约18.7亿欧元，比共识低5\%\par}
{\small\color{refgray}数据：营业利润约30.2亿欧元，比共识高3\%\par}
{\small\color{refgray}边际变化：市场此前关注短期业绩beat，但MS指出成本控制驱动的超预期对估值修复作用有限，重估叙事转向2027年管线催化剂和资本配置战略。\par}
\paragraph{GS} 中国生物制药呼吸管线形成‘自有研发+外部引进’双轮驱动。TQC3721（PDE3/4抑制剂）以2亿美元首付、最高17亿美元里程碑授予阿斯利康中国以外权益，验证平台价值。同步引入GSK两款COPD成熟产品（Trelegy和Anoro），授权期5.5年。TQC3721二期数据：3mg组FEV1改善+100mL，6mg组+147mL，与已上市恩塞芬特林（+124mL）疗效可比，且同步开发干粉吸入剂（DPI）形成差异化。
{\small\color{refgray}数据：TQC3721首付款2亿美元，最高里程碑17亿美元\par}
{\small\color{refgray}数据：二期数据：3mg组FEV1改善+100mL，6mg组+147mL\par}
{\small\color{refgray}数据：恩塞芬特林FEV1改善+124mL\par}
{\small\color{refgray}数据：0例因不良反应停药\par}
{\small\color{refgray}边际变化：此前市场对中国生物制药的认知集中在仿制药集采和管线不确定性，GS指出连续两笔MNC交易（阿斯利康+赛诺菲）正在推动其估值逻辑从‘仿制药企业’向‘创新药平台’切换。\par}
\paragraph{JPM} 脑机接口竞争格局不是技术竞赛，而是临床疗效比拼。行业按患者严重程度、临床需求和设备侵入性分层：非侵入式适合轻度中风康复；半侵入式对严重功能损伤患者具有最佳风险收益比；侵入式（如Neuralink）仅适合小众人群。监管批准不等于商业突破，真正拐点需要形成‘临床数据→算法优化→医生认可→医保覆盖→持续使用’的正向循环。硬件瓶颈在系统集成，而非通道数。
{\small\color{refgray}数据：非侵入式脑电系统信噪比存在根本劣势，但对中轻度中风患者已足够\par}
{\small\color{refgray}数据：半侵入式系统避开硬脑膜穿透，降低手术复杂性\par}
{\small\color{refgray}数据：侵入式系统可及人群始终狭窄\par}
{\small\color{refgray}数据：200通道已够临床用，512通道能覆盖大部分需求\par}
{\small\color{refgray}边际变化：此前市场关注技术指标（如通道数）和监管批准，JPM指出真正的分水岭在于企业能否证明产品在特定适应症上优于现有标准疗法，商业化节奏比近期政策信号所暗示的更慢、更不稳定。\par}
\subsubsection*{关键数据}
\begin{itemize}
  \item 赛诺菲Q2研发费用比共识低5\%，营业利润率27.9\%比共识高90个基点
  \item TQC3721二期FEV1改善+100mL至+147mL，0例因不良反应停药
  \item Trelegy和Anoro 2025年全球合计销售近50亿美元，中国合计不到10亿人民币
  \item 脑机接口200通道已够临床用，512通道能覆盖大部分需求
  \item 中国生物制药呼吸团队超2000人，覆盖超9000家医院
  \item 赛诺菲全年EPS指引可能从‘略高于销售增长’上调至‘低双位数’
\end{itemize}
{\small\color{refgray}本节综合 3 篇研究；涉及 MS、GS、JPM。}
\newpage
\subsection{金融与地产：政策驱动与盈利修复的结构性分化}
\textbf{中国金融与地产板块正经历政策驱动的结构性重估，但内部分化显著：香港银行受益于人民币流动性制度性扩容，内地券商因客户需求驱动的盈利质量提升而获重估，而地产板块的盈利修复弹性被低估；与此同时，美国银行收购借记网络可能重塑支付格局，韩国银行股东回报框架升级成为估值催化剂。}
\subsubsection*{主流共识}
\begin{itemize}
  \item 中国资本外流管控转向“疏堵结合”，资金从灰色渠道转向合规通道，利好创业板和科创板。
  \item 香港人民币流动性安排扩容至5000亿元，是制度性突破，对香港银行构成结构性正面影响。
  \item 中国地产板块盈利修复弹性被低估，毛利率每高出预期0.5个百分点，利润上修13-27\%。
\end{itemize}
\subsubsection*{分歧与条件差异}
\begin{itemize}
  \item 中国银行板块近期上涨是避险轮动还是基本面改善？JPM认为缺乏盈利支撑，行情难以持续；GS则看好香港银行的结构性受益。
  \item 美国银行收购借记网络对Visa/Mastercard的影响程度存在分歧：MS估算单个大行客户流失影响收入0.5-1\%，但网络规模效应难以复制。
  \item 韩国银行股东回报框架升级能否推动估值持续重估？JPM认为当前市净率仍低于1倍，有空间；但市场对息差收窄的担忧仍在。
\end{itemize}
\subsubsection*{机构观点}
\paragraph{GS} 香港人民币流动性安排扩容至5000亿元并新增9个月至3年期品种，是系统性制度安排而非零散利好，对香港银行业构成结构性正面影响。中银香港作为唯一人民币清算行，清算手续费收入将直接受益于业务规模扩大，成为最确定受益者。市场反应验证指向性：中银香港当日收涨5.2\%，其他香港银行平均下跌0.5\%。
{\small\color{refgray}数据：香港金管局人民币业务设施额度从2000亿提升至5000亿元\par}
{\small\color{refgray}数据：新增9个月、2年、3年期品种，此前仅覆盖1个月至1年\par}
{\small\color{refgray}数据：中银香港当日收涨5.2\%，其他香港银行平均下跌0.5\%\par}
{\small\color{refgray}数据：当前香港人民币存款约1万亿元，央行融资额度相当于市场流动性半壁江山\par}
{\small\color{refgray}边际变化：从零散利好认知转向系统性制度安排，中长期流动性供给不足的核心问题得到解决。\par}
\paragraph{HSBC} 中国新一轮资本外流管控是精准引导资金从灰色渠道转向合规通道，对创业板和科创板构成结构性利好。QDII额度持续扩大、南向通资金流入，表明政策意图是“疏堵结合”而非全面收紧。境内财富管理机构是净受益者，因离岸研究活动被迫回流，拥有合规境内渠道的机构将获得增量资金。
{\small\color{refgray}数据：跨境股权TRS暂停对头部上市券商收入影响约0.8-1.2\%，约29-43亿元\par}
{\small\color{refgray}数据：QDII额度与产品净值同步上升\par}
{\small\color{refgray}数据：本轮为2016年以来第三轮资本外流限制周期\par}
{\small\color{refgray}数据：前两轮分别针对房地产/保险（2016-17）和互联网（2021-22）\par}
{\small\color{refgray}边际变化：从简单堵截转向疏堵结合，资金不会凭空消失只会改道，合规通道受益。\par}
\paragraph{HSBC} 中国地产板块盈利修复弹性被市场低估，毛利率每高出预期0.5个百分点，主要开发商2026-2028年盈利可上修13-27\%；高出1个百分点则弹性达26-54\%。政策能见度清晰（低利率环境独立于全球），科技板块财富效应正向高端住宅传导。估值已回落至近期低位，向上超预期空间大于香港地产。
{\small\color{refgray}数据：毛利率每高0.5个百分点，盈利上修13-27\%\par}
{\small\color{refgray}数据：毛利率每高1个百分点，盈利上修26-54\%\par}
{\small\color{refgray}数据：2027年有望迎来广泛盈利复苏\par}
{\small\color{refgray}数据：内地地产板块市净率已回到近期低位\par}
{\small\color{refgray}边际变化：从市场对盈利持续悲观转向关注毛利率改善和减值减轻带来的弹性。\par}
\paragraph{JPM} 中国银行板块近期上涨是地缘避险和AI获利了结驱动的轮动，缺乏盈利支撑，预计将进入横盘整理至8月底业绩期。国有银行二季度营收增速5-8\%、利润增速3-5\%，不足以推动估值持续扩张。渣打回调2\%因中东风险（收入敞口6\%、利润8.1\%）和获利了结，回调提供观察窗口。
{\small\color{refgray}数据：7月8日国有银行H股单日上涨5\%，渣打回调2\%\par}
{\small\color{refgray}数据：过去一个月MSCI中国银行指数跑输亚太银行指数16个百分点\par}
{\small\color{refgray}数据：国有银行二季度营收增速5-8\%，利润增速3-5\%\par}
{\small\color{refgray}数据：地缘事件后一个月中国银行指数平均上涨5\%，跑赢恒生12个百分点\par}
{\small\color{refgray}边际变化：从追涨避险行情转向判断行情不可持续，强调盈利支撑不足。\par}
\paragraph{JPM} CICC二季度净利润同比增长79-102\%至41-46亿元创历史新高，但更关键的是盈利质量——增长由客户需求驱动，而非IPO跟投浮盈，可持续性更强。上半年年化ROE达14-15\%，远超市场全年预期9.8\%，盈利预测存在系统性上修空间。当前0.72倍前瞻市净率较2020-21年牛市1.0-1.3倍仍有差距。
{\small\color{refgray}数据：二季度净利润41-46亿元，同比+79-102\%\par}
{\small\color{refgray}数据：上半年年化ROE 14-15\%，市场预期全年仅9.8\%\par}
{\small\color{refgray}数据：上半年净利润已占JPM全年预测52-55\%\par}
{\small\color{refgray}数据：当前报价0.72倍前瞻市净率，2020-21年牛市为1.0-1.3倍\par}
{\small\color{refgray}边际变化：从关注业绩增速转向关注盈利质量，客户需求驱动模式更具韧性。\par}
\paragraph{JPM} 韩国银行股7月以来表现优于市场8\%，但上涨不是终点而是重估中段。二季度净利润预计约5.9万亿韩元（同比+3\%），息差分化（大行扩大、Kakao Bank反弹6bps），非利息收入走强。真正催化剂是七月密集事件：二季度财报、新一轮回购计划（KB和Shinhan可能更大规模）、韩亚资金价值提升计划。大型金融集团市净率仍低于1倍，ROE 11-12\%有重估空间。
{\small\color{refgray}数据：7月以来韩国银行股上涨8\%，KOSPI下跌15\%\par}
{\small\color{refgray}数据：二季度整体净利润约5.9万亿韩元，同比+3\%\par}
{\small\color{refgray}数据：Kakao Bank息差反弹约6个基点\par}
{\small\color{refgray}数据：大型金融集团市净率仍低于1倍，ROE 11-12\%\par}
{\small\color{refgray}边际变化：从股东回报承诺转向公式化政策，回购规模可能扩大，估值重估进入中段。\par}
\paragraph{MS} 美国大行探索收购Fiserv借记网络STAR和Accel，本质是复制Capital One收购Discover逻辑——通过拥有网络合法收取更高交换费，绕过Durbin修正案管制。Fiserv网络年收入约10亿美元，EBITDA约5亿美元，估值30-40亿美元，溢价20\%以上。真正冲击目标是Chime和Block（Cash App），因大行可将更高交换费收入反哺用户，挤压其核心价值。
{\small\color{refgray}数据：Fiserv借记网络年收入约10亿美元，EBITDA约5亿美元\par}
{\small\color{refgray}数据：资产估值30-40亿美元，按6-8倍EV/EBITDA\par}
{\small\color{refgray}数据：单个大行客户流失对Visa/Mastercard收入影响约0.5-1\%\par}
{\small\color{refgray}数据：借记卡交换费收入占Chime总收入约40\%\par}
{\small\color{refgray}边际变化：从关注Fiserv估值重估转向改写美国借记支付格局，冲击金融科技公司。\par}
\subsubsection*{关键数据}
\begin{itemize}
  \item 香港金管局人民币业务设施额度从2000亿提升至5000亿元
  \item 中银香港当日收涨5.2\%，其他香港银行平均下跌0.5\%
  \item 跨境股权TRS暂停对头部上市券商收入影响约0.8-1.2\%
  \item 毛利率每高0.5个百分点，地产盈利上修13-27\%
  \item CICC上半年年化ROE 14-15\%，市场预期全年仅9.8\%
  \item 7月以来韩国银行股上涨8\%，KOSPI下跌15\%
  \item Fiserv借记网络年收入约10亿美元，EBITDA约5亿美元
  \item 借记卡交换费收入占Chime总收入约40\%
\end{itemize}
\begin{figure}[htbp]
\centering
\includegraphics[width=0.88\linewidth]{figures/fig_065.jpg}
\caption{Exhibit 1 - Exhibit 1: HK system RMB deposit yoy growth looks to be gaining traction As of May'26 data \#\# Disclosure Appendix \#\# Reg AC We, Melissa Kuang, CFA and Wayne Wang, hereby certify that all of the views expressed in this report accu}
\end{figure}
\begin{figure}[htbp]
\centering
\includegraphics[width=0.88\linewidth]{figures/fig_082.jpg}
\caption{Exhibit 1 - Exhibit 1. China's policy measures on outbound investment since 2016 Exhibit 2. Continued southbound inflows during capital outflow restriction periods Exhibit 3. Impact on financial institutions from the Regulation and recent regulatory policies Note: See Hon}
\end{figure}
\begin{figure}[htbp]
\centering
\includegraphics[width=0.88\linewidth]{figures/fig_087.jpg}
\caption{Figure 6 - Figure 6: SOE developers' forward PE (x) Note: The average forward PE multiples of COLI, CRL, Vanke, Greentown, Yuexiu, C\&D Int'l, and Longfor are based on Bloomberg consensus. Figure 7: SOE developers' NAV discount (\%) Figure}
\end{figure}
\begin{figure}[htbp]
\centering
\includegraphics[width=0.88\linewidth]{figures/fig_094.jpg}
\caption{Figure 1 - Figure 1: CICC quarterly net profits reached Rmb4.1-4.6bn in 2Q26, +79-102\% y/y Rmb mn Figure 2: We estimate the annualized ROE has reached \$14 - 15\textbackslash{}\%\$ for CICC in 1H26 Rmb mn Figure 3: CICC's consensus forward P/B compared to}
\end{figure}
{\small\color{refgray}本节综合 7 篇研究；涉及 GS、HSBC、JPM、MS。}
\newpage
\subsection{区域市场：非美与新兴市场的结构性重塑}
\textbf{新兴市场指数已从传统宏观敞口转变为科技集中押注，而日本治理改革正驱动非AI交易量结构性提升，中国市场的反弹则更多源于仓位调整而非外资回流，真正的催化剂需等待7月底至8月。}
\subsubsection*{主流共识}
\begin{itemize}
  \item 新兴市场指数成分已发生根本性变化，前三大重仓股（台积电、三星、SK海力士）占比近30\%，指数集中度达39.4\%，超过标普500的36.7\%。
  \item 日本市场交易活跃度提升不仅由AI驱动，非AI公司换手率从4月175\%升至6月244\%，显示治理改革带来的流动性改善。
  \item 中国近期反弹主要由空头回补和南向资金支撑，而非外资主动回流，全球投资者兴趣上升但行动谨慎。
\end{itemize}
\subsubsection*{分歧与条件差异}
\begin{itemize}
  \item DB认为新兴市场ETF本质上是大型科技押注，而MS强调中国云计算市场正从价格战转向价值竞争，AI云业务成为增长主引擎。
  \item GS认为日本治理改革带来的交易量扩大是可持续的结构性变化，而MS对中国反弹的持续性持谨慎态度，认为需等待7月底至8月的催化剂。
  \item DB指出非美指数长期回报持平，但近期表现相对改善，而MS认为中国市场的反弹缺乏基本面支撑，属于被动减仓结束。
\end{itemize}
\subsubsection*{机构观点}
\paragraph{DB} 新兴市场ETF的底层结构已根本性变化，前三大重仓股（台积电、三星、SK海力士）占指数近30\%权重，前十大集中度达39.4\%，超过标普500的36.7\%。过去近20年，非美全球指数价格与2007年6月持平，而标普500近乎翻四倍。但近期非美市场开始表现相对更好，地缘政治和AI资本支出外溢至电力、硬件和材料等传统宏观板块，为新兴市场提供支撑。
{\small\color{refgray}数据：非美全球指数2025年初价格与2007年6月持平\par}
{\small\color{refgray}数据：标普500同期近乎翻四倍\par}
{\small\color{refgray}数据：新兴市场指数前三大重仓股占比近30\%\par}
{\small\color{refgray}数据：前十大集中度39.4\%，超过标普500的36.7\%\par}
{\small\color{refgray}边际变化：新兴市场ETF从传统宏观敞口转变为科技集中押注，非美市场近期开始表现相对更好。\par}
\paragraph{GS} 日本市场日均交易额4-6月同比增长106\%，剔除AI公司后仍增长49\%-137\%。非AI公司月均换手率从4月175\%升至6月244\%，去年同期仅100\%出头。交易活跃度提升不仅由AI驱动，更反映公司治理改革带来的外资和散户参与增加，是制度性红利带来的流动性改善。
{\small\color{refgray}数据：4-6月日均交易额同比+106\%\par}
{\small\color{refgray}数据：剔除AI后日均交易额同比+49\%至+137\%\par}
{\small\color{refgray}数据：非AI公司月均换手率从4月175\%升至6月244\%\par}
{\small\color{refgray}数据：去年同期非AI换手率仅100\%出头\par}
{\small\color{refgray}边际变化：日本市场交易活跃度从AI主题驱动转向治理改革驱动的结构性改善，非AI公司换手率大幅提升。\par}
\paragraph{MS} 中国近期反弹主要源于空头回补和南向资金支撑，而非外资回流。全球投资者兴趣上升但行动谨慎，预计未来几个月逐步重新配置。真正的催化剂在7月底至8月：电商二季报确认价格战对盈利冲击见顶、AI商业化进展（如腾讯混元3免费AI智能体）、7月IPO解禁压力消化。
{\small\color{refgray}数据：反弹驱动力为空头回补和南向资金\par}
{\small\color{refgray}数据：全球投资者兴趣上升但行动谨慎\par}
{\small\color{refgray}数据：7月底至8月为关键窗口\par}
{\small\color{refgray}数据：电商二季报价格战冲击见顶\par}
{\small\color{refgray}边际变化：中国反弹从外资回流叙事转为仓位调整驱动，催化剂窗口后移至7月底至8月。\par}
\paragraph{MS} 中国云计算市场从价格战转向价值竞争，金山云AI收入高双位数增长，传统云增速降至个位数。AI云业务毛利率高于传统云，长期有望达15\%以上。小米和金山生态提供稳定订单，差异化路径区别于阿里云、腾讯云。行业价格调整（AI计算产品降价15-50\%，存储降价30-50\%）反映算力供不应求。
{\small\color{refgray}数据：AI云收入高双位数增长\par}
{\small\color{refgray}数据：传统云增速降至个位数\par}
{\small\color{refgray}数据：AI云业务毛利率长期有望达15\%以上\par}
{\small\color{refgray}数据：AI计算产品降价15-50\%\par}
{\small\color{refgray}边际变化：中国云计算从价格战转向价值竞争，AI云成为增长主引擎，生态锚定模式提供差异化路径。\par}
\subsubsection*{关键数据}
\begin{itemize}
  \item 非美全球指数2025年初价格与2007年6月持平
  \item 标普500同期近乎翻四倍
  \item 新兴市场指数前十大集中度39.4\%，超过标普500的36.7\%
  \item 日本4-6月日均交易额同比+106\%，剔除AI后仍+49\%至+137\%
  \item 非AI公司月均换手率从4月175\%升至6月244\%
  \item 中国反弹驱动力为空头回补和南向资金
  \item AI云收入高双位数增长，传统云增速降至个位数
  \item AI计算产品降价15-50\%，存储降价30-50\%
\end{itemize}
\begin{figure}[htbp]
\centering
\includegraphics[width=0.88\linewidth]{figures/fig_063.jpg}
\caption{Figure 1 - Figure 1: iShares MSCI ACWI ex U.S. ETF}
\end{figure}
\begin{figure}[htbp]
\centering
\includegraphics[width=0.88\linewidth]{figures/fig_081.jpg}
\caption{Exhibit 3 - Exhibit 3: Monthly turnover for the entire TSE and TSE excluding AI stocks TSE turnover includes ETFs \& REITs Collateral management income Policy rate assumptions are our economist team's}
\end{figure}
\begin{figure}[htbp]
\centering
\includegraphics[width=0.88\linewidth]{figures/fig_064.jpg}
\caption{Figure 1 - Figure 1: iShares MSCI ACWI ex U.S. ETF}
\end{figure}
{\small\color{refgray}本节综合 4 篇研究；涉及 DB、GS、MS。}
\newpage
\subsection{环境与固废：RNG政策套利驱动高利润，信用价格波动成核心变量}
\textbf{美国固废行业的可再生天然气（RNG）业务正从边缘走向核心，其EBITDA利润率高达40-50\%，但90\%的营收依赖D3 RIN信用额度，本质上是政策驱动的套利模式。当前D3 RIN价格约2.70美元/个，投资回收期仅1.5-3年，但EPA年度配额调整和信用价格波动（2-3.75美元区间）将导致EBITDA从5.13亿至10.26亿美元的巨幅变化，构成行业最大不确定性与分歧来源。}
\subsubsection*{主流共识}
\begin{itemize}
  \item RNG业务是固废行业增长最快、利润率最高的细分领域，EBITDA利润率可达40-50\%，远超传统填埋业务。
  \item RNG设施投资回报高度依赖D3 RIN信用价格，当前约2.70美元/个，信用收入占总营收约90\%。
  \item 一个典型RNG设施投资约1.4亿美元（含税收优惠后约1.05亿），在现行信用价格下回收期约1.5-3年。
  \item 每百万英热单位RNG可产生11.727个D3 RIN，信用收入约32.25美元，而天然气商品价值仅约3美元。
\end{itemize}
\subsubsection*{分歧与条件差异}
\begin{itemize}
  \item D3 RIN价格敏感性差异：部分机构认为EPA将维持或提高配额，支撑信用价格在2.70美元以上；另一些机构担忧政策收紧或信用过剩，价格可能跌至2美元以下。
  \item RNG业务规模认知分歧：有观点认为RNG仅占行业营收2-3\%，属于边缘业务；但Bernstein强调其利润贡献和增长潜力远超营收占比，正成为核心利润引擎。
  \item 对政策风险的评估不同：乐观派认为RFS框架稳定，RNG是减排刚需；悲观派指出90\%利润来自政策创设信用，一旦EPA调整规则或信用价格下跌，业务将面临重大冲击。
\end{itemize}
\subsubsection*{机构观点}
\paragraph{Bernstein} 美国固废行业的RNG业务是利润最高、增长最快但波动最大的板块。一个典型RNG设施投资约1.4亿美元（含税收优惠后1.05亿），在D3 RIN价格2.70美元/个下，回收期仅1.5-3年，EBITDA利润率高达40-50\%。但90\%的营收来自D3 RIN信用销售，天然气商品价值仅占10\%。敏感性分析显示，D3 RIN价格从2美元升至3.75美元，EBITDA将从5.13亿美元翻倍至10.26亿美元。当前RNG仅占行业营收2-3\%，但利润贡献和增长潜力远超此比例，正从边缘业务转型为核心利润引擎。
{\small\color{refgray}数据：RNG设施投资约1.4亿美元，含税收优惠后约1.05亿\par}
{\small\color{refgray}数据：D3 RIN当前价格约2.70美元/个\par}
{\small\color{refgray}数据：每MMBtu RNG产生11.727个D3 RIN，信用收入约32.25美元\par}
{\small\color{refgray}数据：天然气商品价值仅约3美元/MMBtu\par}
{\small\color{refgray}边际变化：此前市场将RNG视为固废行业的边缘可再生能源业务，Bernstein首次系统量化其政策套利本质和利润弹性，指出90\%利润来自D3 RIN信用而非天然气销售，并揭示信用价格波动可导致EBITDA翻倍或腰斩，将RNG从成本项重新定义为政策驱动的利润核心。\par}
\subsubsection*{关键数据}
\begin{itemize}
  \item RNG设施投资约1.4亿美元，含税收优惠后约1.05亿
  \item D3 RIN当前价格约2.70美元/个
  \item 每MMBtu RNG产生11.727个D3 RIN，信用收入约32.25美元
  \item 天然气商品价值仅约3美元/MMBtu
  \item EBITDA利润率40-50\%
  \item 投资回收期1.5-3年
  \item D3 RIN价格2-3.75美元区间，EBITDA从5.13亿至10.26亿美元
  \item RNG占行业营收2-3\%
\end{itemize}
\begin{figure}[htbp]
\centering
\includegraphics[width=0.88\linewidth]{figures/fig_011.jpg}
\caption{Exhibit 1 - EXHIBIT 1: The Waste Industry Lifecycle EXHIBIT 2: The Renewable Energy segment is a small but growing part of the industry, driving 2-3\% of industry sales and EBITDA}
\end{figure}
\begin{figure}[htbp]
\centering
\includegraphics[width=0.88\linewidth]{figures/fig_012.jpg}
\caption{Exhibit 2 - EXHIBIT 2: The Renewable Energy segment is a small but growing part of the industry, driving 2-3\% of industry sales and EBITDA Traditional Solid Waste Revenue \& EBITDA Mix Est by Segment EXHIBIT 3: We estimate Renewable Energy ma}
\end{figure}
\begin{figure}[htbp]
\centering
\includegraphics[width=0.88\linewidth]{figures/fig_013.jpg}
\caption{EXHIBIT 3 - EXHIBIT 3: We estimate Renewable Energy margins are attractive at 40-50\%, but subject to changes in pricing of RINs or renewable fuel credits Traditional Solid Waste EBITDA Margin Estimate by Segment}
\end{figure}
\begin{figure}[htbp]
\centering
\includegraphics[width=0.88\linewidth]{figures/fig_014.jpg}
\caption{EXHIBIT 5 - EXHIBIT 7: Methane is captured from landfills and generates a RIN that is sold or it is converted into electricity}
\end{figure}
{\small\color{refgray}本节综合 1 篇研究；涉及 Bernstein。}
\newpage
\subsection{数据中心与基础设施：REIT通道开启与AI资本开支周期}
\textbf{中国数据中心行业正经历双重结构性变化：REIT审批通道的正式开放打通了资产退出闭环，而云厂商AI资本开支的强劲周期则为服务器需求提供了量价支撑。两者共同推动行业从重资产持有向资本效率驱动转型，但受益标的和节奏存在显著分化。}
\subsubsection*{主流共识}
\begin{itemize}
  \item 数据中心REIT审批窗口已正式打开，两单项目获发改委推荐，标志着资产证券化从个案试点进入可复制通道。
  \item 中国云厂商AI资本开支正处于上升周期，预计2026-2028年增速分别为+42\%、+14\%、+10\%，直接拉动AI服务器需求。
  \item 润泽科技通过存量REIT扩募已跑通‘建设-运营-退出’闭环，资产注入年内可期，资本效率优势明显。
\end{itemize}
\subsubsection*{分歧与条件差异}
\begin{itemize}
  \item 万国数据REIT进程时间表不确定性更高，年内完成资产注入的可见度低于润泽科技。
  \item 浪潮信息业绩超预期主要来自‘量’的扩张，但毛利率从2024年6.8\%降至2025年4.9\%，利润率压缩与规模增长并存。
  \item 国内云厂商正从全球级GPU向本地AI芯片过渡，浪潮作为核心受益方，但切换节奏和供应链稳定性仍存变数。
\end{itemize}
\subsubsection*{机构观点}
\paragraph{GS} 中国数据中心REIT审批窗口正式打开，润泽科技廊坊A7、A8数据中心（合计78MW IT容量，利用率近90\%）已纳入存量REIT扩募资产池，计划发行77亿元，估值约80亿元，年内资产注入确定性高。万国数据项目未出现在推荐名单中，参考润泽从获批到上市耗时5-6个月，万国年内完成资产注入的可见度偏低。REIT化打通了‘建设-运营-退出’闭环，但谁能赶上第一波取决于项目成熟度和审批节奏。
{\small\color{refgray}数据：润泽科技廊坊A7、A8数据中心合计IT容量78MW，利用率近90\%\par}
{\small\color{refgray}数据：润泽REIT扩募计划发行金额77亿元，资产估值约80亿元\par}
{\small\color{refgray}数据：润泽从发改委审批到REIT挂牌耗时约5-6个月\par}
{\small\color{refgray}数据：万国数据项目未出现在发改委推荐列表中\par}
{\small\color{refgray}边际变化：REIT审批从个案试点进入可复制通道，但万国数据年内落地概率降低，与管理层近期谨慎表态一致。\par}
\paragraph{GS} 浪潮信息2026年第二季度净利润指引20-25亿元，远超GS此前预期的10亿元，主因中国云厂商AI资本开支强劲。GS预计2026-2028年中国云厂商资本开支增速分别为+42\%、+14\%、+10\%，浪潮作为核心受益方，增长节奏与客户采购周期高度同步。但毛利率从2024年6.8\%降至2025年4.9\%，2026年回升至5.7\%仍远低于2023年10.0\%，规模扩张伴随利润率压缩。存货高企（1Q26达440亿元）为后续交付提供支撑，但也带来不确定性。
{\small\color{refgray}数据：浪潮Q2净利润指引20-25亿元，GS此前预期10亿元\par}
{\small\color{refgray}数据：2026-2028年中国云厂商资本开支增速+42\%/+14\%/+10\%\par}
{\small\color{refgray}数据：毛利率从2024年6.8\%降至2025年4.9\%，2026年回升至5.7\%\par}
{\small\color{refgray}数据：1Q26存货440亿元，4Q25为470亿元\par}
{\small\color{refgray}边际变化：业绩超预期来自AI服务器‘量’的扩张，但毛利率持续低于2023年水平，利润率压缩成为新关注点。\par}
\subsubsection*{关键数据}
\begin{itemize}
  \item 润泽科技REIT扩募计划发行77亿元，资产估值约80亿元
  \item 润泽廊坊A7、A8数据中心合计78MW IT容量，利用率近90\%
  \item 万国数据项目未出现在发改委推荐列表中
  \item 浪潮Q2净利润指引20-25亿元，远超GS预期10亿元
  \item 2026-2028年中国云厂商资本开支增速+42\%/+14\%/+10\%
  \item 浪潮毛利率从2024年6.8\%降至2025年4.9\%
  \item 浪潮1Q26存货440亿元
\end{itemize}
{\small\color{refgray}本节综合 2 篇研究；涉及 GS。}
\newpage
\subsection{面板周期：业绩冲高与价格拐点的博弈}
\textbf{京东方二季度利润超预期主要受非经常性收益驱动，面板主业改善幅度有限；电视面板价格三季度将承压，但供给端纪律提供下行支撑，行业正从业绩高峰向周期拐点过渡。}
\subsubsection*{主流共识}
\begin{itemize}
  \item 电视面板价格从三季度开始将进入调整通道，下半年订单动能低于季节性。
  \item 主要面板厂商通过控制产出维持价格纪律，价格下行空间有限。
  \item 京东方二季度扣非净利润环比增长14-30\%，面板主业处于改善趋势中。
\end{itemize}
\subsubsection*{分歧与条件差异}
\begin{itemize}
  \item 市场对京东方二季度利润超预期的归因存在分歧：部分观点认为归母利润高增反映行业景气，另一部分强调非经常性收益占比近半，主业改善幅度被高估。
  \item 对面板价格拐点后的跌幅判断不一：乐观者认为供给纪律将缓冲下行，谨慎者担忧需求疲软导致跌幅超预期。
\end{itemize}
\subsubsection*{机构观点}
\paragraph{MS} 京东方二季度初步净利润中值35.4亿元，较MS预测高81\%、较市场共识高77\%，但扣非净利润仅16.4-18.7亿元，非经常性收益占比近半，主业改善幅度有限。电视面板价格将从三季度开始进入调整，下半年订单动能低于季节性，面板厂议价能力减弱；不过主要厂商控制产出，价格下行有支撑。公司玻璃基板先进封装布局推进中，计划2027年投资50亿元建设月产能1.5万片产线，目标2028年量产，但收入贡献尚待订单验证。
{\small\color{refgray}数据：京东方2Q26归母净利润预告区间50-55亿元，中值35.4亿元\par}
{\small\color{refgray}数据：扣非净利润16.4-18.7亿元，环比增长14-30\%\par}
{\small\color{refgray}数据：非经常性收益占比约50\%\par}
{\small\color{refgray}数据：MS预测2026/2027/2028年EPS分别为0.20/0.27/0.35元\par}
{\small\color{refgray}边际变化：此前市场聚焦京东方二季度利润超预期，MS报告明确提示价格拐点临近，并拆解利润结构指出非经常性收益主导，边际上从乐观情绪转向对周期拐点和主业真实动能的审视。\par}
\subsubsection*{关键数据}
\begin{itemize}
  \item 京东方2Q26归母净利润中值35.4亿元，较MS预测高81\%、较市场共识高77\%
  \item 扣非净利润16.4-18.7亿元，环比增长14-30\%
  \item 非经常性收益占比约50\%
  \item 电视面板价格三季度开始调整，下半年订单动能低于季节性
  \item 主要面板厂商控制产出，价格下行有支撑
  \item MS预测2026/2027/2028年EPS分别为0.20/0.27/0.35元
  \item ROE从2025年4.4\%提升至2028年8.8\%
  \item 先进封装产线规划月产能1.5万片（510×515mm），目标2028年量产
\end{itemize}
{\small\color{refgray}本节综合 1 篇研究；涉及 MS。}
\newpage
\subsection{小米汽车：插混战略与定价博弈}
\textbf{小米汽车在SU7订单增速放缓的背景下，正式推出插电混动（PHEV）'天行'系列，以家庭和户外场景切入，预计定价20-45万元，入门价或低于20万，直接对标比亚迪豹系列，意图通过PHEV实现走量目标。}
\subsubsection*{主流共识}
\begin{itemize}
  \item 小米'天行'系列定位家庭和户外旅行市场，与SU7/YU7高性能纯电路线形成差异化。
  \item 小米选择PHEV动力形式以缓解长途和户外场景的补能焦虑，降低新品牌信任门槛。
  \item SU7订单热潮消退，小米需依靠'天行'系列走量以实现全年销量目标。
\end{itemize}
\subsubsection*{分歧与条件差异}
\begin{itemize}
  \item 定价策略分歧：部分观点认为'天行'系列入门价可能低于20万以抢占市场，另一观点认为小米可能维持中高端定位以保护品牌形象。
  \item 市场反应差异：乐观预期认为PHEV能快速打开家庭用户市场，谨慎观点则担忧插混市场竞争激烈，小米品牌溢价有限。
\end{itemize}
\subsubsection*{机构观点}
\paragraph{DB} 小米汽车在SU7订单增速放缓（近四周周新增订单从7200降至5400辆）的背景下，正式推出插电混动（PHEV）'天行'系列，首款车型为全尺寸七座SUV'昆仑N90'，后续还有N80和N70。该系列瞄准家庭和户外旅行场景，第二排座椅可180度旋转，车顶可能配备升顶帐篷。DB预计'天行'系列定价在20-45万元区间，入门价可能低于20万，直接对标比亚迪豹系列。小米选择PHEV动力形式以缓解长途补能焦虑，降低新品牌信任门槛，意图通过该系列实现走量目标，预测2026年全年插混交付量可达15万台。电池供应商为欣旺达和中航锂电，成本控制更注重效率。
{\small\color{refgray}数据：小米近四周周新增订单：5000、7200、6100、5400辆\par}
{\small\color{refgray}数据：'天行'系列定价区间：20-45万元\par}
{\small\color{refgray}数据：首款车型'昆仑N90'为全尺寸七座SUV\par}
{\small\color{refgray}数据：2026年插混交付量预测：15万台\par}
{\small\color{refgray}边际变化：从纯电战略转向插电混动，从高性能市场转向家庭和户外场景，定价策略从SU7的20万以上下探至可能低于20万。\par}
\subsubsection*{关键数据}
\begin{itemize}
  \item 小米近四周周新增订单：5000、7200、6100、5400辆
  \item '天行'系列定价区间：20-45万元
  \item 首款车型'昆仑N90'为全尺寸七座SUV
  \item 2026年插混交付量预测：15万台
  \item 电池供应商：欣旺达、中航锂电
\end{itemize}
{\small\color{refgray}本节综合 1 篇研究；涉及 DB。}
\newpage
\subsection{IT服务与软件：宏观不确定性下的结构性分化}
\textbf{宏观不确定性加剧导致IT服务板块整体指引下移，但IBM凭借软件韧性与企业AI需求，成为板块内结构性受益者，与其他依赖可自由支配支出的公司形成分化。}
\subsubsection*{主流共识}
\begin{itemize}
  \item 宏观不确定性自4-5月起显著影响客户决策，IT服务板块整体指引中枢将下移。
  \item 多数公司将主动收窄营收指引上限，市场预期已提前下修。
  \item 企业AI需求是当前板块少数明确的增长驱动力。
\end{itemize}
\subsubsection*{分歧与条件差异}
\begin{itemize}
  \item IBM的软件+咨询组合被视为企业AI需求的净受益者，而EPAM等公司因偏重可自由支配支出而承压。
  \item 指引下修幅度存在信息差：若公司给出高于悲观预期的下限，可能触发修复行情。
  \item 软件业务对AI冲击的防御力强于咨询业务，导致不同子领域表现分化。
\end{itemize}
\subsubsection*{机构观点}
\paragraph{GS} 宏观不确定性自4-5月起明显加剧客户决策周期，IT服务板块整体指引上限将向下调整，市场预期已提前下修（EPAM、Globant、Cognizant近一月分别下跌约12\%、22\%、20\%）。IBM凭借软件韧性（预计2Q软件收入81.56亿美元，全年收入713.25亿美元，固定汇率增长5.2\%）和咨询组合，成为企业AI需求的净受益者，是板块内结构性差异化机会。
{\small\color{refgray}数据：EPAM近一月下跌约12\%\par}
{\small\color{refgray}数据：Globant近一月下跌约22\%\par}
{\small\color{refgray}数据：Cognizant近一月下跌约20\%\par}
{\small\color{refgray}数据：IBM预计2Q软件收入81.56亿美元\par}
{\small\color{refgray}边际变化：从宏观不确定性讨论转向财报季硬约束，指引调整成为核心看点；IBM从板块整体承压中被识别为结构性受益者。\par}
\subsubsection*{关键数据}
\begin{itemize}
  \item EPAM近一月下跌约12\%
  \item Globant近一月下跌约22\%
  \item Cognizant近一月下跌约20\%
  \item IBM预计2Q软件收入81.56亿美元
  \item IBM全年收入713.25亿美元，固定汇率增长5.2\%
  \item IBM目标价335美元，基于25倍远期估值
\end{itemize}
{\small\color{refgray}本节综合 1 篇研究；涉及 GS。}
\newpage
\section{辅助信号｜战略咨询}
本日波士顿咨询报告聚焦三大主题：AI从效率工具向自主运营与战略决策的跃迁、传统零售业的结构性困境与转型路径、以及油气行业估值逻辑的潜在转向。
\subsection{AI进入深水区：从效率工具到自主运营与战略决策}
\textbf{企业AI投入翻倍，但多数资金错配至执行层；真正的结构性价值来自端到端流程重构（Agentic运营）和决策智能体（Decision Agent）进入董事会，后者将重塑企业战略竞争优势。}
\begin{itemize}
  \item 波士顿咨询2026年AI雷达调查显示，企业AI研究占收入比例预计翻倍至1.7\%，但60\%的企业未从AI获得结构性价值，因为流程未做端到端重构，AI收益被锁死在局部。
  \item 传统企业仅不到3\%的工作流程由自主代理运行，而AI原生企业（如按代理架构设计的初创公司）员工与收入比率可达每人每年100万-500万美元，差距源于流程架构而非模型选择。
  \item 企业AI投入存在结构性错配：大部分流向运营和任务自动化，仅不到五分之一的高管将“决策”列为AI首要优先级；在声称“决策是AI战略重点”的企业中，实际投入决策智能体的比例仅14\%。
  \item 决策智能体（Decision Agent）能实时整合数据、模拟情景（如需求走软、供应商提价），将AI从“执行工具”升级为“决策伙伴”，成为下一个竞争壁垒。
  \item 构建决策智能体需要数据层（共享数据湖）、语义层（存储业务逻辑和本体）和安全架构等结构性投入，而非一次性项目。
\end{itemize}
\subsection{传统零售的结构性困境：精选品类零售商以“少”胜“多”}
\textbf{传统超市的扰动是结构性的而非周期性的：消费者转向高频、小批量、就近购买，而精选品类零售商（如Aldi、Lidl）通过系统性的效率碾压（更小面积、更少SKU、更低价格和人工成本）正在夺取市场份额。}
\begin{itemize}
  \item 传统超市占地2500-5000平方米、陈列16000-18000个SKU；精选品类零售商仅需1000-1800平方米、3000-4000个SKU，却能以低20\%-25\%的价格提供同等或更好的商品体验。
  \item 精选品类零售商的人工成本从销售额的9\%-12\%降至4\%-5.5\%，物业成本削减约50\%，优势源于从选址、陈列、供应链到人工配置的系统性效率碾压。
  \item 传统超市的应对——加品类、加促销、砍人工——是三重错误叠加，反而加速了自身边缘化；消费者已用脚投票，选择更快、更便宜、更清晰的购物体验。
  \item 破局关键：建立严格的SKU淘汰机制（而非不断加新品），用“员工有效服务时间”而非“总工时”衡量运营效率，优化不创造价值的后台流程而非一线服务。
  \item 客户中心化需从口号变为组织决策逻辑：BCG跨行业调查发现，只有不到一半的领导层决策真正反映客户洞察，在战略决策中这一比例降至三分之一。
\end{itemize}
\subsection{油气行业估值逻辑的潜在转向：资产组合寿命重新成为焦点}
\textbf{过去十年资本纪律和回报率主导油气公司估值，但“长寿”作为估值驱动因素的权重正在上升，背后是需求韧性、供应管道压缩和地缘政治不确定性的三重叠加。}
\begin{itemize}
  \item BCG报告指出，股东开始重新评估资产组合的寿命，这是自1998-2008年超级周期后首次；尽管变化幅度尚小，但预示着上游市场进入竞争更激烈、战略选择更分化的新阶段。
  \item 过去十年大型油气公司平均五年资本支出与折旧摊销比几乎减半，显著削减了未来产量研究；当全球需求未如预期快速下降，供应管道变薄的影响开始显现。
  \item 地缘政治事件（尤其是霍尔木兹海峡扰动）将能源安全推至前所未有的优先级，进一步强化了“长寿”作为估值驱动因素的重要性。
  \item 过去靠“省钱”获得估值溢价的公司，未来可能需要靠“有选择地花钱”（即投资于长寿命资产组合）来维持地位。
\end{itemize}
\newpage
\section{辅助信号｜智库/国际机构}
本板块整合了亚洲开发银行、布鲁盖尔研究所、IMF、经合组织和世界贸易组织的最新研究，提炼出三个核心主题：全球供应链与贸易格局的结构性重塑、新兴市场与前沿经济体的治理与效率瓶颈、以及绿色转型与AI竞争中的政策与市场错位。这些报告提供了对宏观、产业、企业战略及资产叙事有实质帮助的辅助信号。
\subsection{全球供应链与贸易格局的结构性重塑}
\textbf{全球供应链并未系统性脱钩，而是经历着产业、国家和企业所有权层面的静默重组；同时，贸易规则正在从原则性承诺转向可操作的工具箱，尤其在化石燃料补贴领域。}
\begin{itemize}
  \item 经合组织报告显示，2024年全球价值链贸易占GDP比重（实际值）维持在17\%的历史高位附近，名义值与实际值的背离源于2021-2022年能源价格暴涨后的回落，而非生产网络收缩。
  \item 产业层面分化显著：汽车行业GVC依赖度下降，而服务业和部分高技术制造业的连接度仍在加深；企业并未大规模放弃跨境采购，全球生产网络的物理连接依然牢固。
  \item 世界贸易组织部长级会议将化石燃料补贴改革从联合声明推进至产出文件阶段，核心支柱包括：透明化（贸易政策审议新增标准化提问清单）、扰动应对（设计临时性补贴指南）和识别有害补贴。
  \item 经合组织报告指出，跨国公司海外产出达25万亿美元，全球供应链骨架未动摇，但所有权层面的分析显示，跨国公司在部分战略行业的内部贸易占比上升，反映了企业层面的风险分散策略。
  \item IMF报告揭示，中东和中亚地区的关税是市场权力的加速器：关税越高，企业定价能力增长越快，生产率差距越大；竞争政策、反腐败和产权保护的改善能有效抑制市场权力膨胀。
\end{itemize}
\begin{figure}[htbp]
\centering
\includegraphics[width=0.88\linewidth]{figures/fig_193.jpg}
\caption{Figure 2 - Before the global financial crisis, both nominal and real measures show a strong increase in trade via GVCs, confirming their rapid expansion during the period of “hyperglobalisation” (Subramanian and Kessler, 2014[4]). After 2011, however, the two series dive}
\end{figure}
\begin{figure}[htbp]
\centering
\includegraphics[width=0.88\linewidth]{figures/fig_146.jpg}
\caption{Figure 1 - \#\# 4.4 Stylized Facts \#\# 4.4.1 Market Power in ME\&CA Figure 1. Sales-weighted Markups in ME\&CA and Other EMDEs Sources: Orbis database by Bureau van Dijk, and the authors' estimates. Turning to country-level dynamics, the data reveal meaningful heterogeneity a}
\end{figure}
\subsection{新兴市场与前沿经济体的治理与效率瓶颈}
\textbf{新兴市场和发展中经济体的核心挑战并非资源或资金短缺，而是治理碎片化、执行赤字和资源配置效率低下，这构成了企业战略和资产配置中不可忽视的隐性成本。}
\begin{itemize}
  \item 亚洲开发银行报告指出，喜马拉雅地区住房问题的核心不是短缺，而是抗震与可负担脱节：人口增长正加速向高地震不确定性区域集中，但公共部门未能提供稳定、可预期的成片住宅用地。
  \item IMF对纳米比亚社会支出的分析显示，教育、医疗和社会救助支出水平不低，但效率缺口显著：教育资金锁定在经常性支出，资源向高等教育倾斜，而基础教育生均支出低于新兴市场平均水平。
  \item 经合组织为马耳他制定的技能战略报告指出，核心瓶颈是治理碎片化：决策权分散在多个部委和机构之间，没有一个实体拥有跨部门协调的正式授权，导致数据不共享、需求预测与培训供给脱节。
  \item IMF技术援助报告揭示，伊拉克央行已成功开发宏观预测工具（MFT），但工具尚未进入央行管理层决策会议，也未能与财政部中期预算信息有效对接，处于“实验室”阶段。
  \item 亚洲开发银行研究所报告指出，政府在国际投资争端调解中表现不灵活，核心障碍是问责制与部门利益：官员倾向于严格遵循授权范围，宁可走完耗时耗力的仲裁程序，也不愿在调解中承担“越权”风险。
\end{itemize}
\begin{figure}[htbp]
\centering
\includegraphics[width=0.88\linewidth]{figures/fig_130.jpg}
\caption{Figure 1 - Over the last 3 decades (1992–2022), the Global Emergency Management Database, Emergency Events Database (EM-DAT), recorded 264 disasters in the study area, leading to fatalities (71,691 deaths), displacement (1.2 billion people affected, including 8.4 million}
\end{figure}
\begin{figure}[htbp]
\centering
\includegraphics[width=0.88\linewidth]{figures/fig_151.jpg}
\caption{Figure 1 - Once the forecast discussed and finalized, prepare a draft forecast presentation. \textbackslash{}- Finally, prepare few alternative scenarios which usually focus on different oil price trajectories or fiscal adjustments or any other major macro shock facing Iraqi economy. T}
\end{figure}
\subsection{绿色转型与AI竞争中的政策与市场错位}
\textbf{在绿色转型和AI领域，政策目标与市场现实之间存在显著错位：产地规则可能扭曲资源配置，稳定币和AI收益分配呈现结构性不对称，而开源AI可能并非竞争的解药。}
\begin{itemize}
  \item 布鲁盖尔研究所报告指出，欧盟《产业加速法案》设定的20\%制造业GDP目标缺乏宏观环境逻辑，可能扭曲资源配置；‘欧洲制造’的产地规则将推高成本、延迟脱碳，并面临WTO挑战。
  \item IMF报告显示，AI每年节省2.7万亿美元劳动成本，但发展中国家几乎未分到：低收入国家AI对话集中在极窄职业范围，无法像高收入国家那样通过大规模就业基数放大收益。
  \item 经合组织报告揭示，AI模型层竞争加剧（开发者从9家增至47家，前沿玩家从2家扩至8家），但上游芯片、算力和数据资源高度集中，下游竞争随时可能逆转；开源AI可能成为另一种垄断陷阱。
  \item IMF报告指出，稳定币在新兴市场是双刃剑：当汇率失衡较低时，稳定币改善外汇准入和价格发现；当失衡较高时，公共价格信号会压缩信念分散度，放大挤兑风险。
  \item 经合组织报告显示，绿色债券（GSSS）的投资者兴趣真实，但实际配置仍集中在发达市场：2025年发展中国家仅占全球绿色债券的17\%（剔除中国后降至5\%），流动性是投资者最在意的隐性门槛。
\end{itemize}
\begin{figure}[htbp]
\centering
\includegraphics[width=0.88\linewidth]{figures/fig_134.jpg}
\caption{Figure 1 - Figure 1: European battery cell manufacturing capacity by company, region of HQ and operational status (as of Q1 2026) FDI screening has historically been a national competence. The EU introduced only in 2019 coordination of scr}
\end{figure}
\begin{figure}[htbp]
\centering
\includegraphics[width=0.88\linewidth]{figures/fig_179.jpg}
\caption{Figure 1 - Empirical evidence suggests continuous innovation and declining prices, especially in the AI model development segment. Despite significant entry costs and the presence of a few tech giants in the AI model development segment \$\textasciicircum{}\{1\}\$ , the performance of Genera}
\end{figure}
\subsection{宏观政策与金融稳定}
\textbf{用于校准利率、通胀、财政约束和金融稳定叙事。}
\begin{itemize}
  \item 一份来自亚洲开发银行（ADB）的专题简报，没有聚焦宏观预测或利率走势，而是花大量篇幅拆解了一个看似技术性的基础设施——Bond Connect。但读完你会发现，它讨论的远不止是一个交易通道。
  \item 一份来自IMF技术援助团队的最新评估报告，对多米尼加共和国的财政透明度做了全面诊断。结论是：这个国家的财政报告和预算流程已有相当基础，但真正的不确定性藏在报表之外——公共信托产品、长期购电协议、养老金缺口和国企财务状况，这些关键领域要么数据缺失，要么分析流于形式。报告的核心判断是...
  \item 一份2025年5月完成、2026年1月发布的IMF技术援助报告，揭示了一个多数市场观察者容易忽略的事实：在非洲南部小国史瓦帝尼，非银行机构的资产规模已占据整个资金体系的66.2\%，但此前长期游离在官方货币与资金统计框架之外。报告的核心判断是，IMF正在帮助该国央行将养老金、保险、...
  \item 伊拉克宏观环境长期被一个问题困扰：高度依赖石油收入，油价波动直接冲击财政、货币和国际收支。中央银行需要一套能预测这些冲击、支持政策决策的分析框架。IMF的技术援助报告给出的结论很直接——工具已经有了，真正的挑战不在模型本身，而在团队能否持续用它、数据能否支撑它、央行能否与财政部真...
  \item 在固定汇率制度下，货币政策独立性天然受到制约。但这份IMF最新工作论文揭示了一个更具体的结论：纳米比亚的银行存贷款利率变动，主导力量来自南非储备银行的政策利率，而非本国央行的回购利率。这不仅是学术验证，更对任何实行货币局制度或紧密盯住汇率的地区，提供了一个关于政策传导边界的现实参...
  \item 一份来自国际货币产品组织（IMF）的最新工作论文，系统拆解了美国回购市场利差的驱动因素。结论并不复杂，但含义值得细品：准备金依然是回购市场最核心的稳定器，但它的效果并非在所有市场环境下都相同。真正决定回购利率走向的，是准备金、交易商资产负债表使用率、对冲产品杠杆与国债发行这四个变...
  \item 2025年3月，联合国统计委员会与国际货币产品组织同步发布了国民账户体系和国际收支手册的最新版本。对大多数国家的统计机构而言，这不仅仅是一次技术升级，而是一场需要跨部门协调、用户参与和资源重新配置的系统工程。IMF统计部门最新发布的《SNA/BPM实施规划手册》给出了一个清晰的信...
  \item 理解资金体系之外的宏观环境活动，对政策制定者而言从来不是可选项。但加密货币挖矿——这个不依赖传统资金中介、直接创造新资产的入口——其真实规模在全球范围内几乎是一团迷雾。IMF最新工作论文提出了一种全新的测量思路：通过追踪全球主要产地的矿机出口数据，来反推各国的挖矿活动分布。这个方...
\end{itemize}
\begin{figure}[htbp]
\centering
\includegraphics[width=0.88\linewidth]{figures/fig_127.jpg}
\caption{Figure 3 - Figure 3: Breakdown of Northbound Bond Connect Investors by Jurisdiction Figure 4: Breakdown of Northbound Bond Connect Investors by Investor Type Figure 5: Assets Under Custody by International Institutional Investors in the}
\end{figure}
\begin{figure}[htbp]
\centering
\includegraphics[width=0.88\linewidth]{figures/fig_156.jpg}
\caption{Figure 1 - \#\# 5. The close tracking of the BoN's policy rate with that of the SARB has changed in recent years. Prior to 2021, the BoN repo rate closely tracked the SARB policy rate and was typically set at a modest premium, reflecting the BoN's assessment of domestic \#\#}
\end{figure}
\subsection{地缘政治与安全}
\textbf{用于补充能源、航运、防务和供应链风险的尾部情景。}
\begin{itemize}
  \item 全球水安全研究的回报率已被反复证明——联合国估算，每投入1美元，平均能产生约4美元的收益。但现实是，全球水相关基础设施长期面临资金不足、资产过早老化、服务交付效率低下等问题。经合组织（OECD）最新发布的《水安全融资》报告给出了一个核心判断：问题不在于缺钱，而在于钱怎么投、怎么回...
  \item 一份来自经合组织的最新跟进报告揭示了土耳其反海外贿赂体系中的一个尖锐矛盾：政府持续做出改革承诺，但实际执行进展却几近停滞。在2024年第四轮评估提出的71项建议中，土耳其仅完全落实了10项，仍有49项未采取任何实质性行动。这是该工作组自2007年以来反复敦促却始终未见突破的领域。
\end{itemize}
\begin{figure}[htbp]
\centering
\includegraphics[width=0.88\linewidth]{figures/fig_183.jpg}
\caption{Figure 2 - Unlike GSS bonds, which allocate proceeds to specific projects, “sustainability-linked” bonds (SLBs) are general-purpose bonds, whose financial terms are tied to the achievement of environmental and social performance targets (OECD, 2024[9]). In other words, t}
\end{figure}
\section{结语}
后续需验证的数据与叙事节点包括：美国中期选举后政策不确定性回落对市场风格切换的持续性；亚洲厄尔尼诺对农村消费的实际冲击幅度；韩元弱势是否因美元整体走弱而缓解；铝市场2027年过剩预期是否因中东产能进一步加速而扩大；人形机器人量产节奏是否如期在2026年Q3达到每周1000台；EPA年度配额调整对D3 RIN信用价格的影响；以及中国云计算市场从价格战转向价值竞争的实际进展。
\newpage
\section{覆盖概览}
投行/券商 59 篇；战略咨询 5 篇；智库/国际机构 23 篇。\par
投行覆盖：GS 14篇 / MS 13篇 / JPM 7篇 / Citi 6篇 / Bernstein 5篇 / DB 5篇 / NOM 5篇 / HSBC 2篇 / BARC 1篇 / UBS 1篇\par
\section{Disclaimer}
{\scriptsize\color{refgray}
This community is a paid learning and information-sharing community created for general educational purposes only. The membership fee is charged solely for access to community discussions, educational content, organization, and operational costs. It is not an advisory fee, management fee, performance fee, brokerage fee, or compensation for personalized financial, investment, legal, tax, or professional advice.\par
I participate and share information solely as an individual and for educational discussion among community members. I am not acting as a financial adviser, investment adviser, broker, dealer, portfolio manager, legal adviser, tax adviser, accountant, fiduciary, or any other licensed professional adviser. Nothing shared in this community constitutes, and should not be interpreted as, financial advice, investment advice, legal advice, tax advice, a recommendation, solicitation, offer, endorsement, instruction, or guarantee to buy, sell, hold, short, trade, or invest in any security, cryptocurrency, fund, derivative, strategy, or other financial product.\par
All content shared in this community, including messages, discussions, links, files, charts, examples, market commentary, watchlists, AI-generated or AI-polished summaries, and third-party materials, is general, impersonal, and educational in nature. It does not take into account any individual's financial situation, investment objectives, risk tolerance, investment horizon, liquidity needs, tax status, legal restrictions, or suitability requirements. No content should be relied upon as suitable for any specific person, account, portfolio, or financial decision.\par
Information shared in this community may be incomplete, inaccurate, outdated, speculative, AI-generated, AI-edited, or based on third-party internet sources that have not been independently verified. I make no representation or warranty as to the accuracy, completeness, timeliness, reliability, or suitability of any information shared.\par
Financial markets involve substantial risk, including the possible loss of principal. Past performance, historical data, hypothetical examples, simulated results, backtests, personal opinions, or educational examples do not guarantee future results. Each member is solely responsible for conducting their own research, exercising independent judgment, and making their own decisions. Before making any financial, investment, legal, or tax decision, members should consult qualified and licensed professionals.\par
Participation in this community, payment of any membership fee, or communication with me or other members does not create any adviser-client, fiduciary, professional, contractual, agency, partnership, employment, or other advisory relationship. By joining or participating in this community, you acknowledge and agree that you are solely responsible for your own decisions, actions, transactions, profits, losses, risks, and consequences, and that I shall not be liable for any direct, indirect, incidental, consequential, financial, legal, tax, or other loss or damage arising from or related to any content, discussion, information, or material shared in this community.\par
}
\newpage
\begin{center}
{\Large 更多详情报告kcdesk.com}\par\vspace{0.8cm}
\includegraphics[width=0.58\linewidth]{\detokenize{../../prompts/zsxq_img.jpg}}
\end{center}
\end{document}
